%%%%%%%%%%%%%%%%%%%%%%%%%%%%%%%%%%%%%%%%%%%%%%%%%%%%%%%%%%%%%%%%%%%%%%%%%%%%%%%%%%
\begin{frame}[fragile]\frametitle{}
\begin{center}
{\Large Demonstration Skills}
\end{center}
\end{frame}

%%%%%%%%%%%%%%%%%%%%%%%%%%%%%%%%%%%%%%%%%%%%%%%%%%%%%%%%%%%
\begin{frame}[fragile]\frametitle{Syllabus}

\begin{itemize}
\item 4.1  Prayer: Concept and recitation of \textit{Pranava} (\textit{प्रणव}) and hymns. 
\item 4.2  Yoga Cleansing Techniques: Knowledge of Dhauti (\textit{धौती}), Neti (\textit{नेति}) and practice of Kapalabhati (\textit{कपालभाति}). 
\item 4.3  Yogic Sukshma Vyayama (\textit{सुक्ष्म व्यायाम}) and Sthula Vyayama (\textit{स्थूल व्यायाम})      
	\begin{itemize}
	\item a. Yogic Sukshma Vyayama (Micro Circulation Practices) 
		\begin{itemize}
		\item Neck Movement: Griva Shakti Vikasaka (\textit{ग्रीवा शक्ति विकासक}) (I, II, III, IV) 
		\item Shoulder Movement: Bhuja Valli Shakti Vikasaka (\textit{भुजा वल्ली शक्ति विकासक}), Purna Bhuja Shakti Vikasaka (\textit{पूर्ण भुजा शक्ति विकासक}) 
		\item Trunk Movement: Kati Shakti Vikasaka (\textit{कटि शक्ति विकासक}) (I, II, III, IV, V) 
		\item Knee Movement: Jangha Shakti Vikasaka (\textit{जङ्घा शक्ति विकासक}) (II-A \& B), Janu Shakti Vikasaka (\textit{जानु शक्ति विकासक})
		\item Ankle Movement: Pada-mula Shakti Vikasaka (\textit{पाद-मूल शक्ति विकासक}) - A \& B, Gulpha-pada-pristha-pada-tala Shakti Vikasaka (\textit{गुल्फ-पाद-पृष्ठ-पाद-तल शक्ति विकासक})
		\end{itemize}
	\item b. Yogic Sthula Vyayama (Macro Circulation Practices) 
		\begin{itemize}
		\item Sarvanga Pushti (\textit{सर्वाङ्ग पुष्टि})
		\item Hrid Gati (\textit{हृद् गति}) (Engine Daud)
		\end{itemize}
	\end{itemize}
\end{itemize}
\end{frame}


%%%%%%%%%%%%%%%%%%%%%%%%%%%%%%%%%%%%%%%%%%%%%%%%%%%%%%%%%%%
\begin{frame}[fragile]\frametitle{Syllabus}

\begin{itemize}
\item 4.4  Yogic Surya Namaskara (\textit{योगिक सूर्य नमस्कार})   
\item 4.5  Yogasana (\textit{योगासन}) 
	\begin{itemize}
	\item Tadasana (\textit{ताड़ासन}), Vrikshasana (\textit{वृक्षासन}), Ardha Chakrasana (\textit{अर्ध चक्रासन}), Padahastasana (\textit{पादहस्तासन}), Kati Chakrasana (\textit{कटिचक्रासन}), Trikonasana (\textit{त्रिकोणासन})  
	\item Dandasana (\textit{दंडासन}), Sukhasana (\textit{सुखासन}), Padmasana (\textit{पद्मासन}), Vajrasana (\textit{वज्रासन})
	\item Bhadrasana (\textit{भद्रासन}), Mandukasana (\textit{मण्डुकासन}), Ushtrasana (\textit{उष्ट्रासन}), Shashankasana (\textit{शशांकासन}), Uttana Mandukasana (\textit{उत्तान मण्डुकासन}) 
	\item Paschimottanasana (\textit{पश्चिमोत्तानासन}), Purvottanasana (\textit{पूर्वोत्तानासन}) 
	\item Vakrasana (\textit{वक्रासन}), Gomukhasana (\textit{गोमुखासन}) 
	\item Bhujangasana (\textit{भुजंगासन}), Shalabhasana (\textit{शलभासन}), Makarasana (\textit{मकरासन}) 
	\item Pavanamuktasana (\textit{पवनमुक्तासन}), Uttanapadasana (\textit{उत्तानपादासन}), Ardha Halasana (\textit{अर्ध हलासन}), Setubandhasana (\textit{सेतुबन्धासन}) 
	\item Vipareetakarani (\textit{विपरीतकरणी}), Saral Matsyasana (\textit{सरल मत्स्यासन}), Shavasana (\textit{शवासन})  
	\end{itemize}

\item 4.6  Preparatory Breathing Practices: Sectional breathing (abdominal, thoracic, and clavicular), Yogic deep breathing
	% \begin{itemize}
	% \item Sectional breathing (abdominal, thoracic, and clavicular) 
	% \item Yogic deep breathing 
	% \end{itemize}
	
\item 4.7  Pranayama (\textit{प्राणायाम}) 
	\begin{itemize}
	\item Concept of Puraka (\textit{पूरक}), Rechaka (\textit{रेचक}), and Kumbhaka (\textit{कुम्भक}) 
	\item Anuloma Viloma/Nadi Shodhana (\textit{अनुलोम विलोम/नाड़ी शोधन}), Sheetali (\textit{शीतली}) (without Kumbhaka) Bhramari (\textit{भ्रामरी}) (without Kumbhaka) 
	\end{itemize}
	
\end{itemize}
\end{frame}

%%%%%%%%%%%%%%%%%%%%%%%%%%%%%%%%%%%%%%%%%%%%%%%%%%%%%%%%%%%
\begin{frame}[fragile]\frametitle{Syllabus}

\begin{itemize}
\item 4.8  Understanding of Bandha (\textit{बन्ध})
	\begin{itemize}
	\item Jalandhara Bandha (\textit{जलंधर बन्ध})
	\item Uddiyana Bandha (\textit{उड्डियान बन्ध})
	\item Mula Bandha (\textit{मूल बन्ध})
	\end{itemize}

\item 4.9  Understanding of Mudra (\textit{मुद्रा})   
	\begin{itemize}
	\item Hasta Mudras (\textit{हस्त मुद्राः}) (Chin (\textit{चिन्}), Chinmaya (\textit{चिन्मय}), Brahma (\textit{ब्रह्म}), Adi (\textit{आदि}), Jnana (\textit{ज्ञान}), Dhyana (\textit{ध्यान}), and Nasika (\textit{नासिका}))
	\end{itemize}
	
\item 4.10  Practices leading to Meditation and Dhyana Sadhana (\textit{ध्यान साधना})   
	\begin{itemize}
	\item  Recitation of Pranava (\textit{प्रणव}) \& Soham (\textit{सोऽहम्}) 
	\item  Recitation of selected hymns, invocations, and prayers from Vedas (\textit{वेदाः}) \& Upanishads (\textit{उपनिषदः})
	\item  Body and breath awareness  
	\item  Yoga Nidra (\textit{योग निद्रा})
	\end{itemize}
	
\end{itemize}
\end{frame}


%%%%%%%%%%%%%%%%%%%%%%%%%%%%%%%%%%%%%%%%%%%%%%%%%%%%%%%%%%%%%%%%%%%%%%%%%%%%%%%%%%
\begin{frame}[fragile]\frametitle{}
\begin{center}
{\Large  Prayer: Concept and recitation of प्रणव Pranava and hymns}
\end{center}
\end{frame}


%%%%%%%%%%%%%%%%%%%%%%%%%%%%%%%%%%%%%%%%%%%%%%%%%%%%%%%%%%%
\begin{frame}[fragile]\frametitle{Concept of प्रणव Pranava}\footnotesize
\begin{columns}
    \begin{column}[T]{0.6\linewidth}
      \begin{itemize}
        \item \textbf{प्रणव Pranava}: the syllable \textbf{Om} (ॐ), also written \textbf{AUM}.
        \item Patanjali names \textbf{Pranava} as the word for \textbf{Ishvara}; the Guidebook follows him.
        \item \textbf{AUM}: in the Mandukya Upanishad (9 to 12) \textbf{A} is waking, \textbf{U} dreaming, \textbf{M} deep sleep, and the silence after it the fourth state.
        \item The syllabus lists mantra and japa among the \textbf{practices leading to meditation}.
        \item \textbf{Mantra Japa}: repeated chanting of \textbf{Om} (ॐ), done with its meaning in mind.
        \item Traditionally said to remove obstacles and turn the mind \textbf{inward}.
        \item Traditionally said to steady the mind and bring \textbf{calmness}.
        \item Used in \textbf{meditation} and \textbf{prayer}.
        \item In the tradition, it deepens the \textbf{connection} with the divine.
      \end{itemize}
    \end{column}
    \begin{column}[T]{0.4\linewidth}
        \begin{center}
        % FIG ycb1-fig-pranav | now: images/ycb1_pranav.jpg | final: books/YCB1/images/new/ycb1_pranav_new.pdf | register: books/YCB1/docs/image_register.csv
        % PROMPT SP-DIAGRAM (common text: books/YCB1/docs/image_system_prompt.md) | Used in block ycb1-pd-prayer-pranava-hymns, frame 'Concept of Pranava'. Concept diagram: Concept of Pranava. Boxes, circles and arrows only, drawn from the structure given in the notes block (the names and their order are listed there); Sanskrit terms in Roman letters with the English meaning in brackets, nothing else added. Flat colours that still read in grayscale. Best redrawn later in TikZ. Aspect ratio 4:3. Check before use: author and the Sanskrit check in the notes.
        % NEW-PIC approved by the author (2026-11-27); the original is commented out below it
        \includegraphics[width=\linewidth,keepaspectratio]{new/ycb1_pranav_new.pdf}
        % original, replaced (author approved 2026-11-27): \includegraphics[width=\linewidth,keepaspectratio]{ycb1_pranav}

		% credit of the replaced original: {\tiny (Ref: Certification  of Yoga Professionals Official Guidebook For Level I (Instructor))}
		{\tiny (Line diagram drawn for this deck.)}
        \end{center}	
    \end{column}
\end{columns}
\end{frame}

%%%%%%%%%%%%%%%%%%%%%%%%%%%%%%%%%%%%%%%%%%%%%%%%%%%%%%%%%%%
\begin{frame}[fragile]\frametitle{Concept of Hymns}
\begin{columns}
    \begin{column}[T]{0.5\linewidth}
      \begin{itemize}
        \item \textbf{Hymns}: sacred \textbf{verses} used in \textbf{ritual} and \textbf{devotional} practice.
        \item The syllabus asks you to recite selected hymns, invocations and prayers from the \textbf{Vedas} and \textbf{Upanishads}.
        \item Say the meaning first, then chant slowly and let the class repeat.
        % \item \textbf{Vedic Hymns}: Ancient, transcendental \textbf{sound}. % no source in the notes
        % \item Traditionally used for \textbf{devotion} and \textbf{blessings}. % dropped in the notes (pass 4): no Board or Guidebook source
        \item Keep the prayer inclusive: a call for harmony and well-being.
        \item The Common Yoga Protocol opens with a \textbf{prayer} (a Rigveda hymn) and closes with a \textbf{sankalpa} and the \textbf{shanti path}.
      \end{itemize}
    \end{column}
    \begin{column}[T]{0.5\linewidth}
        \begin{center}
        % FIG ycb1-fig-hymns | now: images/ycb1_hymns.jpg | final: books/YCB1/images/new/ycb1_hymns_new.pdf | register: books/YCB1/docs/image_register.csv
        % PROMPT SP-DIAGRAM (common text: books/YCB1/docs/image_system_prompt.md) | Used in block ycb1-pd-prayer-pranava-hymns, frame 'Concept of Hymns'. Concept diagram: Concept of Hymns. Boxes, circles and arrows only, drawn from the structure given in the notes block (the names and their order are listed there); Sanskrit terms in Roman letters with the English meaning in brackets, nothing else added. Flat colours that still read in grayscale. Best redrawn later in TikZ. Aspect ratio 4:3. Check before use: author and the Sanskrit check in the notes.
        % NEW-PIC (uncomment to show the new picture beside the original; comment the original only after review) \includegraphics[width=0.35\linewidth,keepaspectratio]{new/ycb1_hymns_new.pdf}
        \includegraphics[width=\linewidth,keepaspectratio]{ycb1_hymns}

		{\tiny (Ref: Certification  of Yoga Professionals Official Guidebook For Level I (Instructor))}	
        \end{center}	
    \end{column}
\end{columns}
\end{frame}

% %%%%%%%%%%%%%%%%%%%%%%%%%%%%%%%%%%%%%%%%%%%%%%%%%%%%%%%%%%%
% \begin{frame}[fragile]\frametitle{Recitation of Hymns}
% \begin{columns}
    % \begin{column}[T]{0.5\linewidth}
      % \begin{itemize}
        % \item \textbf{Sanskrit Hymns}: Pronunciation and \textbf{intonation}.
        % \item Enhances \textbf{mental clarity} and \textbf{focus}.
        % \item \textbf{Spiritual growth} through regular recitation.
        % \item Commonly chanted during \textbf{Puja} and \textbf{Satsang}.
        % \item \textbf{Harmonizes} mind and spirit.
      % \end{itemize}
    % \end{column}
    % \begin{column}[T]{0.5\linewidth}
        % \begin{center}
        % \includegraphics[width=\linewidth,keepaspectratio]{ycb1_recite}

		% {\tiny (Ref: Certification  of Yoga Professionals Official Guidebook For Level I (Instructor))}	
        % \end{center}	
    % \end{column}
% \end{columns}
% \end{frame}


%%%%%%%%%%%%%%%%%%%%%%%%%%%%%%%%%%%%%%%%%%%%%%%%%%%%%%%%%%%%%%%%%%%%%%%%%%%%%%%%%%
\begin{frame}[fragile]\frametitle{}
\begin{center}
{\Large Yoga Cleansing Techniques}
\end{center}
\end{frame}

%%%%%%%%%%%%%%%%%%%%%%%%%%%%%%%%%%%%%%%%%%%%%%%%%%%%%%%%%%%
\begin{frame}[fragile]\frametitle{धौती Dhauti}
\begin{columns}
    \begin{column}[T]{0.5\linewidth}
      \begin{itemize}
        \item \textbf{Dhauti} (\textit{धौती}): cleansing of the \textbf{digestive tract} up to the stomach.
        \item One of the six cleansing actions (shatkarma).
        \item Learned only from a \textbf{qualified teacher}; the exam asks for its meaning and types, not a demonstration.
        \item Never taught in a class for beginners.
        \item Traditionally used for \textbf{cleansing}, not as a cure. Three types:
		\begin{itemize}
		  \item \textbf{Vaman Dhauti} (\textit{वमन धौती}): lukewarm salted water.
		  \item \textbf{Danda Dhauti} (\textit{दण्ड धौती}): a soft tube (the Gheranda Samhita names a plant stalk).
		  \item \textbf{Vastra Dhauti} (\textit{वस्त्र धौती}): a moist strip of cloth, swallowed and drawn out under a teacher.
		\end{itemize}		
      \end{itemize}
    \end{column}
    \begin{column}[T]{0.5\linewidth}
        \begin{center}
        % FIG ycb1-fig-dhauti | now: images/ycb1_dhauti.jpg | final: books/YCB1/images/new/ycb1_dhauti_new.png | register: books/YCB1/docs/image_register.csv
        % PROMPT SP-DIAGRAM (common text: books/YCB1/docs/image_system_prompt.md) | Used in block ycb1-pd-cleansing-techniques, frame 'Dhauti'. Cleansing technique: Dhauti. A clean schematic only: the practitioner (CS1) in the starting posture of the technique with a small simple inset showing the path or the area involved, as a flat outline with no internal organs drawn in detail and nothing graphic. Show the posture and the direction of the action with thin arrows. No tools or fluids drawn realistically. Aspect ratio 4:3. This technique is taught only under a qualified teacher; the picture is a memory aid. Check before use: certified teacher confirms the form against the steps in the notes block.
        % NEW-PIC (uncomment to show the new picture beside the original; comment the original only after review) \includegraphics[width=0.35\linewidth,keepaspectratio]{new/ycb1_dhauti_new.png}
        \includegraphics[width=\linewidth,keepaspectratio]{ycb1_dhauti}
		
		{\tiny (Ref: What is Shatkarma? 6 Types of Shatkarma for Purification and Their Benefits - Yogi Anurag)}	
        \end{center}	
    \end{column}
\end{columns}
\end{frame}


%%%%%%%%%%%%%%%%%%%%%%%%%%%%%%%%%%%%%%%%%%%%%%%%%%%%%%%%%%%
\begin{frame}[fragile]\frametitle{नेति Neti}
\begin{columns}
    \begin{column}[T]{0.5\linewidth}
      \begin{itemize}
          \item \textbf{Neti} (\textit{नेति}) Kriya cleanses the nasal passages; learned only from a qualified teacher.
          \item Two types of Neti:
            \begin{itemize}
              \item \textbf{Jala Neti} (\textit{जल नेति}): Uses a neti pot with tepid salt water, poured through one nostril and out the other.
              \item \textbf{Sutra Neti} (\textit{सूत्र नेति}): Uses a smooth thread or rubber tube passed through the nose. Kapalabhati follows neti to clear leftover water.
            \end{itemize}
      \end{itemize}
    \end{column}
    \begin{column}[T]{0.5\linewidth}
        \begin{center}
        % FIG ycb1-fig-neti | now: images/ycb1_neti.png | final: books/YCB1/images/new/ycb1_neti_new.png | register: books/YCB1/docs/image_register.csv
        % PROMPT SP-DIAGRAM (common text: books/YCB1/docs/image_system_prompt.md) | Used in block ycb1-pd-cleansing-techniques, frame 'Neti'. Cleansing technique: Neti. A clean schematic only: the practitioner (CS1) in the starting posture of the technique with a small simple inset showing the path or the area involved, as a flat outline with no internal organs drawn in detail and nothing graphic. Show the posture and the direction of the action with thin arrows. No tools or fluids drawn realistically. Aspect ratio 4:3. This technique is taught only under a qualified teacher; the picture is a memory aid. Check before use: certified teacher confirms the form against the steps in the notes block.
        % NEW-PIC (uncomment to show the new picture beside the original; comment the original only after review) \includegraphics[width=0.35\linewidth,keepaspectratio]{new/ycb1_neti_new.png}
        \includegraphics[width=\linewidth,keepaspectratio]{ycb1_neti}
		
		{\tiny (Ref: What is Shatkarma? 6 Types of Shatkarma for Purification and Their Benefits - Yogi Anurag)}	
        \end{center}	
    \end{column}
\end{columns}
\end{frame}


%%%%%%%%%%%%%%%%%%%%%%%%%%%%%%%%%%%%%%%%%%%%%%%%%%%%%%%%%%%
\begin{frame}[fragile,shrink=6]\frametitle{कपालभाती Kapalabhati}
\begin{columns}
    \begin{column}[T]{0.5\linewidth}
      \begin{itemize}
          \item \textbf{Kapalabhati} (\textit{कपालभाती}): kapala is skull, bhati is shining, so ``skull shining''.
          \item Cleanses the frontal region of the head and the lungs (the Common Yoga Protocol says the frontal air sinuses).
          \item A cleansing kriya that is also taught as a breathing practice (\textit{कपालभाती प्राणायाम}).
          \item Sit in a meditative posture. Breathe out forcefully by pulling the abdomen in; let the in-breath come by itself.
          \item Normal breathing has an active inhalation; Kapalabhati reverses it, with an active exhalation and a passive inhalation.
          \item Common Yoga Protocol: thirty rapid breaths, then one deep breath; three rounds (beginners up to three rounds of twenty).
          \item \textbf{Avoid} with heart conditions, giddiness, high blood pressure, vertigo, chronic nose bleeding, epilepsy, migraine, stroke, hernia or gastric ulcer.
      \end{itemize}
    \end{column}
    \begin{column}[T]{0.5\linewidth}
        \begin{center}
        % FIG ycb1-fig-kapalbhati | now: images/ycb1_kapalbhati.png | final: books/YCB1/images/new/ycb1_kapalbhati_new.png | register: books/YCB1/docs/image_register.csv
        % PROMPT SP-FIG (common text: books/YCB1/docs/image_system_prompt.md) | Used in block ycb1-pd-cleansing-techniques, frame 'Kapalabhati (Skull Shining Breath); Kapalabhati'. Practice: Kapalabhati (Skull Shining Breath); Kapalabhati. One adult practitioner (CS1), seated, front three-quarter view, in the posture of the notes block. Show the key action in a small inset: the abdomen drawn in sharply for kapalabhati, the throat narrowed softly for ujjayi, the steady gaze on a candle flame at eye level for trataka, the tongue position for the mukh mudra. Use only the inset that matches Kapalabhati (Skull Shining Breath); Kapalabhati. Aspect ratio 4:5 portrait. Check before use: certified teacher confirms the form against the steps in the notes block.
        % NEW-PIC (uncomment to show the new picture beside the original; comment the original only after review) \includegraphics[width=0.35\linewidth,keepaspectratio]{new/ycb1_kapalbhati_new.png}
        \includegraphics[width=0.8\linewidth,keepaspectratio]{ycb1_kapalbhati}
		
		{\tiny (Ref: What is Shatkarma? 6 Types of Shatkarma for Purification and Their Benefits - Yogi Anurag)}	
        \end{center}	
    \end{column}
\end{columns}
\end{frame}


%%%%%%%%%%%%%%%%%%%%%%%%%%%%%%%%%%%%%%%%%%%%%%%%%%%%%%%%%%%
\begin{frame}[fragile]\frametitle{नौली Nauli}\footnotesize
\begin{columns}
    \begin{column}[T]{0.6\linewidth}
      \begin{itemize}
          \item \textbf{Nauli} (\textit{नौली}) Kriya: churning and massage of the abdomen.
          \item Learned only from a qualified teacher; never taught in a class for beginners.
          \item The Guidebook gives four forms: madhya, vama, dakshina and nauli chalana (rotation).
          \item The middle abdominal muscles, the \textbf{rectus abdominis} (\textit{अन्तरत्वक}), are separated and isolated.
          \item The muscles are drawn to the middle, the left or the right of the linea alba (\textit{लिनिया अल्बा}).
          \item Basic types of Nauli:
            \begin{itemize}
              \item \textbf{Madhya Nauli} (\textit{मध्य नौली}): muscles drawn together at the centre (linea alba).
              \item \textbf{Vama Nauli} (\textit{वाम नौली}): muscles drawn to the left of the centre.
              \item \textbf{Dakshina Nauli} (\textit{दक्षिण नौली}): muscles drawn to the right of the centre.
            \end{itemize}
      \end{itemize}
    \end{column}
    \begin{column}[T]{0.4\linewidth}
        \begin{center}
        % FIG ycb1-fig-nauli | now: images/ycb1_nauli.png | final: books/YCB1/images/new/ycb1_nauli_new.png | register: books/YCB1/docs/image_register.csv
        % PROMPT SP-DIAGRAM (common text: books/YCB1/docs/image_system_prompt.md) | Used in block ycb1-pd-cleansing-techniques, frame 'Nauli'. Cleansing technique: Nauli. A clean schematic only: the practitioner (CS1) in the starting posture of the technique with a small simple inset showing the path or the area involved, as a flat outline with no internal organs drawn in detail and nothing graphic. Show the posture and the direction of the action with thin arrows. No tools or fluids drawn realistically. Aspect ratio 4:3. This technique is taught only under a qualified teacher; the picture is a memory aid. Check before use: certified teacher confirms the form against the steps in the notes block.
        % NEW-PIC (uncomment to show the new picture beside the original; comment the original only after review) \includegraphics[width=0.35\linewidth,keepaspectratio]{new/ycb1_nauli_new.png}
        \includegraphics[width=\linewidth,keepaspectratio]{ycb1_nauli}
		
		{\tiny (Ref: What is Shatkarma? 6 Types of Shatkarma for Purification and Their Benefits - Yogi Anurag)}	
        \end{center}	
    \end{column}
\end{columns}	  
\end{frame}


%%%%%%%%%%%%%%%%%%%%%%%%%%%%%%%%%%%%%%%%%%%%%%%%%%%%%%%%%%%
\begin{frame}[fragile]\frametitle{त्राटक Trataka}\footnotesize
\begin{columns}
    \begin{column}[T]{0.5\linewidth}
     \begin{itemize}
          \item \textbf{Trataka} (\textit{त्राटक}) Kriya: traditionally used for the eyes and for concentration.
          \item Gaze steadily at one point without blinking.
          \item Gaze at a small mark or a flame, traditionally till the eyes water; a candle is common today. Epileptic persons should not gaze at a flame.
          \item Learned from a qualified teacher.
          \item Two types of Trataka (a common teaching split, not in the Guidebook):
            \begin{itemize}
              \item \textbf{Internal Trataka} (\textit{आन्तरिक त्राटक}): Focus on \textbf{trikuti} (\textit{त्रिकुटी}) (third eye) between eyebrows.
              \item \textbf{External Trataka} (\textit{बाह्य त्राटक}): Gazing at an external point such as a candle flame or a small mark.
            \end{itemize}
      \end{itemize}
    \end{column}
    \begin{column}[T]{0.5\linewidth}
        \begin{center}
        % FIG ycb1-fig-tratak | now: images/ycb1_tratak.png | final: books/YCB1/images/new/ycb1_tratak_new.png | register: books/YCB1/docs/image_register.csv
        % PROMPT SP-FIG (common text: books/YCB1/docs/image_system_prompt.md) | Used in block ycb1-pd-cleansing-techniques, frame 'Trataka'. Practice: Trataka. One adult practitioner (CS1), seated, front three-quarter view, in the posture of the notes block. Show the key action in a small inset: the abdomen drawn in sharply for kapalabhati, the throat narrowed softly for ujjayi, the steady gaze on a candle flame at eye level for trataka, the tongue position for the mukh mudra. Use only the inset that matches Trataka. Aspect ratio 4:5 portrait. Check before use: certified teacher confirms the form against the steps in the notes block.
        % NEW-PIC (uncomment to show the new picture beside the original; comment the original only after review) \includegraphics[width=0.35\linewidth,keepaspectratio]{new/ycb1_tratak_new.png}
        \includegraphics[width=\linewidth,keepaspectratio]{ycb1_tratak}
		
		{\tiny (Ref: What is Shatkarma? 6 Types of Shatkarma for Purification and Their Benefits - Yogi Anurag)}	
        \end{center}	
    \end{column}
\end{columns}	  
\end{frame}

%%%%%%%%%%%%%%%%%%%%%%%%%%%%%%%%%%%%%%%%%%%%%%%%%%%%%%%%%%%
\begin{frame}[fragile]\frametitle{बस्ती Basti}
\begin{columns}
    \begin{column}[T]{0.5\linewidth}
      \begin{itemize}
          \item \textbf{Basti} (\textit{बस्ती}) Kriya cleanses the large intestine; learned only from a qualified teacher, never in a beginners' class.
          \item Two types of Basti:
            \begin{itemize}
              \item \textbf{Sthala Basti} (\textit{स्थल बस्ती}): uses air
              \item \textbf{Jala Basti} (\textit{जल बस्ती}): uses water
            \end{itemize}
          \item Water (jala) or air (sthala) is drawn into the large intestine.
          \item The abdomen is churned while it is held inside.
          \item It is then expelled.
          \item Traditionally used for cleansing the colon.
      \end{itemize}
    \end{column}
    \begin{column}[T]{0.5\linewidth}
        \begin{center}
        % FIG ycb1-fig-basti | now: images/ycb1_basti.jpg | final: books/YCB1/images/new/ycb1_basti_new.png | register: books/YCB1/docs/image_register.csv
        % PROMPT SP-DIAGRAM (common text: books/YCB1/docs/image_system_prompt.md) | Used in block ycb1-pd-cleansing-techniques, frame 'Basti'. Cleansing technique: Basti. A clean schematic only: the practitioner (CS1) in the starting posture of the technique with a small simple inset showing the path or the area involved, as a flat outline with no internal organs drawn in detail and nothing graphic. Show the posture and the direction of the action with thin arrows. No tools or fluids drawn realistically. Aspect ratio 4:3. This technique is taught only under a qualified teacher; the picture is a memory aid. Check before use: certified teacher confirms the form against the steps in the notes block.
        % NEW-PIC (uncomment to show the new picture beside the original; comment the original only after review) \includegraphics[width=0.35\linewidth,keepaspectratio]{new/ycb1_basti_new.png}
        \includegraphics[width=\linewidth,keepaspectratio]{ycb1_basti}
		
		{\tiny (Ref: What is Shatkarma? 6 Types of Shatkarma for Purification and Their Benefits - Yogi Anurag)}	
        \end{center}	
    \end{column}
\end{columns}	  
\end{frame}


%%%%%%%%%%%%%%%%%%%%%%%%%%%%%%%%%%%%%%%%%%%%%%%%%%%%%%%%%%%%%%%%%%%%%%%%%%%%%%%%%%
\begin{frame}[fragile]\frametitle{}
\begin{center}
{\Large योगिक सूक्ष्म व्यायाम  Yogic Sukshma Vyayama and स्थूल व्यायाम Sthula Vyayama}
\end{center}
\end{frame}


%%%%%%%%%%%%%%%%%%%%%%%%%%%%%%%%%%%%%%%%%%%%%%%%%%%%%%%%%%%
\begin{frame}[fragile]\frametitle{Sukshma Vyayama: Concept}
\begin{columns}
    \begin{column}[T]{0.5\linewidth}
      \begin{itemize}
        \item \textbf{Sukshma Vyayama}: subtle \textbf{exercise}; small, slow movements of the \textbf{joints}.
        \item Also called loosening or micro-circulation practices.
        \item Traditionally credited with \textbf{joint mobility}.
        \item Prepares the body for \textbf{stronger practice}.
        \item Used as a \textbf{warm-up} at the start of a session.
		  \begin{itemize}
			\item \textbf{Neck}: Greeva Shakti Vikasaka series.
			\item \textbf{Shoulder}: Bhuja Valli and Purna Bhuja Shakti Vikasaka.
			\item \textbf{Trunk}: Kati Shakti Vikasaka series.
			\item \textbf{Knee and thigh}: Janu and Jangha Shakti Vikasaka.
			\item \textbf{Ankle and foot}: Pada-mula and Gulpha-pada-pristha-pada-tala Shakti Vikasaka.
		  \end{itemize}		
      \end{itemize}
    \end{column}
    \begin{column}[T]{0.5\linewidth}
        \begin{center}
			% FIG ycb1-fig-sukshma | now: images/ycb1_sukshma.jpg | final: books/YCB1/images/new/ycb1_sukshma_new.png | register: books/YCB1/docs/image_register.csv
			% PROMPT SP-FIG (common text: books/YCB1/docs/image_system_prompt.md) | Used in block ycb1-pd-sukshma-vyayama, frame 'Sukshma Vyayama: Concept'. Concept illustration: Sukshma Vyayama: Concept. One practitioner (CS1) performing the practice described in the notes block, plus one small clean schematic if the concept needs it, nothing else. Aspect ratio 4:3. Check before use: certified teacher confirms the form against the steps in the notes block.
			% NEW-PIC (uncomment to show the new picture beside the original; comment the original only after review) \includegraphics[width=0.35\linewidth,keepaspectratio]{new/ycb1_sukshma_new.png}
			\includegraphics[width=\linewidth,keepaspectratio]{ycb1_sukshma}
				
			{\tiny (Ref: Sukshma vyayama: The 7-minute relaxation exercise Activating the Joints)}		
        \end{center}	
    \end{column}
\end{columns}
\end{frame}

%%%%%%%%%%%%%%%%%%%%%%%%%%%%%%%%%%%%%%%%%%%%%%%%%%%%%%%%%%%
\begin{frame}[fragile]\frametitle{Neck Movement}\footnotesize
\begin{columns}
    \begin{column}[T]{0.5\linewidth}
      \begin{itemize}
		\item \textbf{Greeva Shakti Vikasaka I (ग्रीवा शक्ति विकासक I)}: bend forward and back; exhale down, inhale up.
		\item \textbf{Greeva Shakti Vikasaka II (ग्रीवा शक्ति विकासक II)}: twist right and left; exhale to turn, inhale back.
		\item \textbf{Greeva Shakti Vikasaka III (ग्रीवा शक्ति विकासक III)}: tilt right and left; exhale to tilt, inhale back.
		\item \textbf{Greeva Shakti Vikasaka IV (ग्रीवा शक्ति विकासक IV)}: rotate the neck slowly.
		\item \textbf{Note}: two rounds each, slow, with the breath, shoulders relaxed. The Common Yoga Protocol order is forward and back, side bend, twist, rotation; numerals vary between sources.
		\item \textbf{Caution}: gentle with neck pain; a chair may be used. Elderly people and those with cervical problems or high blood pressure should avoid it.
	  \end{itemize}
	  
				{\tiny (Ref: Day 02 of 30 Days of Yogic Journey - Guiding Principles for Yoga Practitioners and Yogic Sukshma Vyayama - Saatvik Life)}		  
    \end{column}
    \begin{column}[T]{0.5\linewidth}
		\begin{center}
		I, II, IV
		
		        % FIG ycb1-fig-greeva | now: images/ycb1_greeva.jpg | final: books/YCB1/images/new/ycb1_greeva_new.png | register: books/YCB1/docs/image_register.csv
		        % PROMPT SP-FIG (common text: books/YCB1/docs/image_system_prompt.md) | Used in block ycb1-pd-sukshma-vyayama, frame 'Neck Movement'. Movement practice: Neck Movement. One adult practitioner (CS1), standing, front or side view as the joint needs. Show the start position and the end position of the movement in one picture, as two overlapping figures in light and normal tone, with one thin curved arrow for the direction. The moving joint is clearly visible; the rest of the body stays still. No props. Aspect ratio 3:2 landscape. Check before use: certified teacher confirms the form against the steps in the notes block.
		        % NEW-PIC (uncomment to show the new picture beside the original; comment the original only after review) \includegraphics[width=0.35\linewidth,keepaspectratio]{new/ycb1_greeva_new.png}
		        \includegraphics[width=0.5\linewidth,keepaspectratio]{ycb1_greeva}
				
				
		III
		
				% FIG ycb1-fig-greevabend | now: images/ycb1_greevabend.jpg | final: books/YCB1/images/new/ycb1_greevabend_new.png | register: books/YCB1/docs/image_register.csv
				% PROMPT SP-FIG (common text: books/YCB1/docs/image_system_prompt.md) | Used in block ycb1-pd-sukshma-vyayama, frame 'Neck Movement'. Movement practice: Neck Movement. One adult practitioner (CS1), standing, front or side view as the joint needs. Show the start position and the end position of the movement in one picture, as two overlapping figures in light and normal tone, with one thin curved arrow for the direction. The moving joint is clearly visible; the rest of the body stays still. No props. Aspect ratio 3:2 landscape. Check before use: certified teacher confirms the form against the steps in the notes block.
				% NEW-PIC (uncomment to show the new picture beside the original; comment the original only after review) \includegraphics[width=0.35\linewidth,keepaspectratio]{new/ycb1_greevabend_new.png}
				\includegraphics[width=0.5\linewidth,keepaspectratio]{ycb1_greevabend}

	
		\end{center}	
    \end{column}
\end{columns}

\end{frame}


%%%%%%%%%%%%%%%%%%%%%%%%%%%%%%%%%%%%%%%%%%%%%%%%%%%%%%%%%%%
\begin{frame}[fragile]\frametitle{Shoulder Movement}\footnotesize
\begin{columns}
    \begin{column}[T]{0.5\linewidth}
      \begin{itemize}
		\item \textbf{Bhuja Valli Shakti Vikasaka (भुजा वल्ली शक्ति विकासक)}: Right arm, then left arm, raised sideways overhead and lowered, palm open and outward, fingers together
		\item \textbf{Bhuja Valli Shakti Vikasaka (भुजा वल्ली शक्ति विकासक)}: Both arms together the same way; arms must not touch the head or each other
		\item \textbf{Purna Bhuja Shakti Vikasaka (पूर्ण भुजा शक्ति विकासक)}: Fists with thumbs tucked in; swing one arm in a full circle forward, then reverse, then the other arm
		\item \textbf{Purna Bhuja Shakti Vikasaka (पूर्ण भुजा शक्ति विकासक)}: Both fists together describe a full circle upward, then reverse
		\item \textbf{Benefits (लाभ)}: Loosens the shoulders and arms (Ref: Brahmachari, Yogic Sukshma Vyayama)
	  \end{itemize}
    \end{column}
    \begin{column}[T]{0.5\linewidth}
		\begin{center}
		        % FIG ycb1-fig-bhuja | now: images/ycb1_bhuja.png | final: books/YCB1/images/new/ycb1_bhuja_new.png | register: books/YCB1/docs/image_register.csv
		        % PROMPT SP-FIG (common text: books/YCB1/docs/image_system_prompt.md) | Used in block ycb1-pd-sukshma-vyayama, frame 'Shoulder Movement'. Movement practice: Shoulder Movement. One adult practitioner (CS1), standing, front or side view as the joint needs. Show the start position and the end position of the movement in one picture, as two overlapping figures in light and normal tone, with one thin curved arrow for the direction. The moving joint is clearly visible; the rest of the body stays still. No props. Aspect ratio 3:2 landscape. Check before use: certified teacher confirms the form against the steps in the notes block.
		        % NEW-PIC (uncomment to show the new picture beside the original; comment the original only after review) \includegraphics[width=0.35\linewidth,keepaspectratio]{new/ycb1_bhuja_new.png}
		        \includegraphics[width=0.7\linewidth,keepaspectratio]{ycb1_bhuja}
				
				{\tiny (Ref: Day 02 of 30 Days of Yogic Journey - Guiding Principles for Yoga Practitioners and Yogic Sukshma Vyayama - Saatvik Life)}		
		\end{center}	
    \end{column}
\end{columns}
\end{frame}


%%%%%%%%%%%%%%%%%%%%%%%%%%%%%%%%%%%%%%%%%%%%%%%%%%%%%%%%%%%
\begin{frame}[fragile,shrink=6]\frametitle{Trunk Movement}\footnotesize
\begin{columns}
    \begin{column}[T]{0.5\linewidth}
      \begin{itemize}
		\item \textbf{Kati Shakti Vikasaka I (कटी शक्ति विकासक I)}: Fist behind the back, other hand on the wrist; bend back with the inhale, then forward with the exhale, head towards the knees.
		\item \textbf{Kati Shakti Vikasaka II (कटी शक्ति विकासक II)}: Legs wide apart, hands on the hips; bend back with the inhale, forward with the exhale, head towards the floor.
		\item \textbf{Kati Shakti Vikasaka III (कटी शक्ति विकासक III)}: Feet together, backward and forward bends without the hands touching the thighs or knees; begin gently.
		\item \textbf{Kati Shakti Vikasaka IV (कटी शक्ति विकासक IV)}: Arms spread sideways; side bends to the left and right, reaching towards the calf.
		\item \textbf{Kati Shakti Vikasaka V (कटी शक्ति विकासक V)}: Feet two feet apart; swing the trunk and outstretched arms in a semicircle to each side.
		\item \textbf{Benefits (लाभ)}: traditionally credited with a mobile spine and waist. \textbf{Caution}: for trunk practice the Common Yoga Protocol says avoid with severe back pain, disc disorders, after abdominal surgery and during menstruation.
	  \end{itemize}
    \end{column}
    \begin{column}[T]{0.5\linewidth}
		\begin{center}
		        % FIG ycb1-fig-kati | now: images/ycb1_kati.png | final: books/YCB1/images/new/ycb1_kati_new.png | register: books/YCB1/docs/image_register.csv
		        % PROMPT SP-FIG (common text: books/YCB1/docs/image_system_prompt.md) | Used in block ycb1-pd-sukshma-vyayama, frame 'Trunk Movement'. Movement practice: Trunk Movement. One adult practitioner (CS1), standing, front or side view as the joint needs. Show the start position and the end position of the movement in one picture, as two overlapping figures in light and normal tone, with one thin curved arrow for the direction. The moving joint is clearly visible; the rest of the body stays still. No props. Aspect ratio 3:2 landscape. Check before use: certified teacher confirms the form against the steps in the notes block.
		        % NEW-PIC (uncomment to show the new picture beside the original; comment the original only after review) \includegraphics[width=0.35\linewidth,keepaspectratio]{new/ycb1_kati_new.png}
		        \includegraphics[width=0.5\linewidth,keepaspectratio]{ycb1_kati}
				
				{\tiny (Ref: SUKSHMA VYAYAMA. STHULA VYAYAMA. YOGA. WARM UP. MICRO CIRCULATION. MACRO CIRCULATION. YOGIC WARM UP. - Living 150)}	
		\end{center}	
    \end{column}
\end{columns}
\end{frame}

%%%%%%%%%%%%%%%%%%%%%%%%%%%%%%%%%%%%%%%%%%%%%%%%%%%%%%%%%%%
\begin{frame}[fragile]\frametitle{Knee Movement}
\begin{columns}
    \begin{column}[T]{0.5\linewidth}
      \begin{itemize}
		\item \textbf{Janu Shakti Vikasaka (जनु शक्ति विकासक)}: Standing, swing the foot forward with a knee jerk, then back so the heel touches the buttock; both legs, upper body erect
		\item \textbf{Janu Shakti Vikasaka (जनु शक्ति विकासक)}: Ten times to begin with
		\item \textbf{Benefits (लाभ)}: Loosens the knee joints. The Common Yoga Protocol knee squat: avoid in acute arthritis.
	  \end{itemize}
	  
	  		{\tiny (Ref: Day 03 of 30 Days of Yogic Journey - Guiding Principles for Yoga Practitioners and Yogic Sukshma Vyayama - Saatvik Life)}	
    \end{column}
    \begin{column}[T]{0.5\linewidth}
		\begin{center}
		% FIG ycb1-fig-janu-1 | now: images/ycb1_janu_1.png | final: books/YCB1/images/new/ycb1_janu_1_new.png | register: books/YCB1/docs/image_register.csv
		% PROMPT SP-FIG (common text: books/YCB1/docs/image_system_prompt.md) | Used in block ycb1-pd-sukshma-vyayama, frame 'Knee Movement'. Movement practice: Knee Movement. One adult practitioner (CS1), standing, front or side view as the joint needs. Show the start position and the end position of the movement in one picture, as two overlapping figures in light and normal tone, with one thin curved arrow for the direction. The moving joint is clearly visible; the rest of the body stays still. No props. Aspect ratio 3:2 landscape. Check before use: certified teacher confirms the form against the steps in the notes block.
		% NEW-PIC (uncomment to show the new picture beside the original; comment the original only after review) \includegraphics[width=0.35\linewidth,keepaspectratio]{new/ycb1_janu_1_new.png}
		\includegraphics[width=0.4\linewidth,keepaspectratio]{ycb1_janu_1}
		
		% FIG ycb1-fig-janu-2 | now: images/ycb1_janu_2.png | final: books/YCB1/images/new/ycb1_janu_2_new.png | register: books/YCB1/docs/image_register.csv
		% PROMPT SP-FIG (common text: books/YCB1/docs/image_system_prompt.md) | Used in block ycb1-pd-sukshma-vyayama, frame 'Knee Movement'. Concept illustration: Knee Movement. One practitioner (CS1) performing the practice described in the notes block, plus one small clean schematic if the concept needs it, nothing else. Aspect ratio 4:3. Check before use: certified teacher confirms the form against the steps in the notes block.
		% NEW-PIC (uncomment to show the new picture beside the original; comment the original only after review) \includegraphics[width=0.35\linewidth,keepaspectratio]{new/ycb1_janu_2_new.png}
		\includegraphics[width=0.4\linewidth,keepaspectratio]{ycb1_janu_2}

		\end{center}	
    \end{column}
\end{columns}
\end{frame}

%%%%%%%%%%%%%%%%%%%%%%%%%%%%%%%%%%%%%%%%%%%%%%%%%%%%%%%%%%%
\begin{frame}[fragile]\frametitle{Thigh Movement}\footnotesize
\begin{columns}
    \begin{column}[T]{0.5\linewidth}
      \begin{itemize}
		\item \textbf{Jangha Shakti Vikasaka I (जंघ शक्ति विकासक I)}: Jumping jacks (in Brahmachari's text; the Board syllabus names only the II practices)
		\item \textbf{Jangha Shakti Vikasaka II-A (जंघ शक्ति विकासक II-A)}: Feet together, bend the knees with the arms held forward until the thighs are parallel to the floor; heels stay down
		\item \textbf{Jangha Shakti Vikasaka II-B (जंघ शक्ति विकासक II-B)}: Heels raised, arms spread sideways, bend the knees apart without sitting on the heels
		\item \textbf{Benefits (लाभ)}: Loosens the thigh and knee joints
	  \end{itemize}
	  
		{\tiny (Ref: Day 03 of 30 Days of Yogic Journey - Guiding Principles for Yoga Practitioners and Yogic Sukshma Vyayama - Saatvik Life)}		  
    \end{column}
    \begin{column}[T]{0.5\linewidth}
		\begin{center}
		% FIG ycb1-fig-jangha-1 | now: images/ycb1_jangha_1.png | final: books/YCB1/images/new/ycb1_jangha_1_new.png | register: books/YCB1/docs/image_register.csv
		% PROMPT SP-FIG (common text: books/YCB1/docs/image_system_prompt.md) | Used in block ycb1-pd-sukshma-vyayama, frame 'Thighs Movement'. Movement practice: Thighs Movement. One adult practitioner (CS1), standing, front or side view as the joint needs. Show the start position and the end position of the movement in one picture, as two overlapping figures in light and normal tone, with one thin curved arrow for the direction. The moving joint is clearly visible; the rest of the body stays still. No props. Aspect ratio 3:2 landscape. Check before use: certified teacher confirms the form against the steps in the notes block.
		% NEW-PIC (uncomment to show the new picture beside the original; comment the original only after review) \includegraphics[width=0.35\linewidth,keepaspectratio]{new/ycb1_jangha_1_new.png}
		\includegraphics[width=0.5\linewidth,keepaspectratio]{ycb1_jangha_1}
		
		% FIG ycb1-fig-jangha-2 | now: images/ycb1_jangha_2.png | final: books/YCB1/images/new/ycb1_jangha_2_new.png | register: books/YCB1/docs/image_register.csv
		% PROMPT SP-FIG (common text: books/YCB1/docs/image_system_prompt.md) | Used in block ycb1-pd-sukshma-vyayama, frame 'Thighs Movement'. Movement practice: Thighs Movement. One adult practitioner (CS1), standing, front or side view as the joint needs. Show the start position and the end position of the movement in one picture, as two overlapping figures in light and normal tone, with one thin curved arrow for the direction. The moving joint is clearly visible; the rest of the body stays still. No props. Aspect ratio 3:2 landscape. Check before use: certified teacher confirms the form against the steps in the notes block.
		% NEW-PIC (uncomment to show the new picture beside the original; comment the original only after review) \includegraphics[width=0.35\linewidth,keepaspectratio]{new/ycb1_jangha_2_new.png}
		\includegraphics[width=0.5\linewidth,keepaspectratio]{ycb1_jangha_2}


		\end{center}	
    \end{column}
\end{columns}
\end{frame}


%%%%%%%%%%%%%%%%%%%%%%%%%%%%%%%%%%%%%%%%%%%%%%%%%%%%%%%%%%%
\begin{frame}[fragile]\frametitle{Ankle Movement}
 \begin{columns}
    \begin{column}[T]{0.5\linewidth}
      \begin{itemize}
		\item \textbf{Pada-mula Shakti Vikasaka A (पाद-मूल शक्ति विकासक A)}: Rise on the toes and lower again, spring-like
		\item \textbf{Pada-mula Shakti Vikasaka B (पाद-मूल शक्ति विकासक B)}: Jump on toes % 
		\item \textbf{Gulpha-pada-pristha-pada-tala Shakti Vikasaka (गुल्फ-पाद-पृष्ठ-पाद-तल शक्ति विकासक)}: Foot held a little above the ground, circle the ankle each way; then the other foot
		% \item \textbf{Gulpha-pada-pristha-pada-tala Shakti Vikasaka (गुल्फ-पाद-पृष्ठ-पाद-तल शक्ति विकासक)}: Heel raises
		\item \textbf{Benefits (लाभ)}: Loosens the ankle and foot joints
	  \end{itemize}
    \end{column}
    \begin{column}[T]{0.5\linewidth}
		\begin{center}
		% FIG ycb1-fig-pada | now: images/ycb1_pada.png | final: books/YCB1/images/new/ycb1_pada_new.png | register: books/YCB1/docs/image_register.csv
		% PROMPT SP-FIG (common text: books/YCB1/docs/image_system_prompt.md) | Used in block ycb1-pd-sukshma-vyayama, frame 'Ankle Movement'. Movement practice: Ankle Movement. One adult practitioner (CS1), standing, front or side view as the joint needs. Show the start position and the end position of the movement in one picture, as two overlapping figures in light and normal tone, with one thin curved arrow for the direction. The moving joint is clearly visible; the rest of the body stays still. No props. Aspect ratio 3:2 landscape. Check before use: certified teacher confirms the form against the steps in the notes block.
		% NEW-PIC (uncomment to show the new picture beside the original; comment the original only after review) \includegraphics[width=0.35\linewidth,keepaspectratio]{new/ycb1_pada_new.png}
		\includegraphics[width=0.5\linewidth,keepaspectratio]{ycb1_pada}
		
		
		% FIG ycb1-fig-gulpha | now: images/ycb1_gulpha.png | final: books/YCB1/images/new/ycb1_gulpha_new.png | register: books/YCB1/docs/image_register.csv
		% PROMPT SP-FIG (common text: books/YCB1/docs/image_system_prompt.md) | Used in block ycb1-pd-sukshma-vyayama, frame 'Ankle Movement'. Movement practice: Ankle Movement. One adult practitioner (CS1), standing, front or side view as the joint needs. Show the start position and the end position of the movement in one picture, as two overlapping figures in light and normal tone, with one thin curved arrow for the direction. The moving joint is clearly visible; the rest of the body stays still. No props. Aspect ratio 3:2 landscape. Check before use: certified teacher confirms the form against the steps in the notes block.
		% NEW-PIC (uncomment to show the new picture beside the original; comment the original only after review) \includegraphics[width=0.35\linewidth,keepaspectratio]{new/ycb1_gulpha_new.png}
		\includegraphics[width=0.5\linewidth,keepaspectratio]{ycb1_gulpha}
		
		\end{center}	
    \end{column}
\end{columns}
\end{frame}


%%%%%%%%%%%%%%%%%%%%%%%%%%%%%%%%%%%%%%%%%%%%%%%%%%%%%%%%%%%
\begin{frame}[fragile]\frametitle{Sthula Vyayama: Concept}
\begin{columns}
    \begin{column}[T]{0.5\linewidth}
      \begin{itemize}
        \item \textbf{Sthula Vyayama}: Gross \textbf{exercise} in Yoga.
        \item Larger loosening movements of the \textbf{gross body}.
        \item The Board syllabus names two practices: Sarvanga Pushti and Hrid Gati.
        \item Strength and stamina are traditionally credited aims.
        \item More strenuous than the gentle sukshma vyayama.
			  \begin{itemize}
				\item \textbf{Sarvanga Pushti}: whole-body practice, next frame.
				\item \textbf{Hrid Gati}: running-like practice, next frame.
				
				
				
			  \end{itemize}		
      \end{itemize}
    \end{column}
    \begin{column}[T]{0.5\linewidth}
        \begin{center}
		% FIG ycb1-fig-sthula | now: images/ycb1_sthula.png | final: books/YCB1/images/new/ycb1_sthula_new.png | register: books/YCB1/docs/image_register.csv
		% PROMPT SP-FIG (common text: books/YCB1/docs/image_system_prompt.md) | Used in block ycb1-pd-sthula-vyayama, frame 'Sthula Vyayama: Concept'. Concept illustration: Sthula Vyayama: Concept. One practitioner (CS1) performing the practice described in the notes block, plus one small clean schematic if the concept needs it, nothing else. Aspect ratio 4:3. Check before use: certified teacher confirms the form against the steps in the notes block.
		% NEW-PIC (uncomment to show the new picture beside the original; comment the original only after review) \includegraphics[width=0.35\linewidth,keepaspectratio]{new/ycb1_sthula_new.png}
		\includegraphics[width=\linewidth,keepaspectratio]{ycb1_sthula}
		
		{\tiny (Ref: Day 06 of 30 Days of Yogic Journey - Guiding Principles for Yoga Practitioners and Yogic Sukshma Vyayama - Saatvik Life)}	
        \end{center}	
    \end{column}
\end{columns}
\end{frame}

% %%%%%%%%%%%%%%%%%%%%%%%%%%%%%%%%%%%%%%%%%%%%%%%%%%%%%%%%%%%
% \begin{frame}[fragile]\frametitle{Examples of Sthula Vyayama}
% \begin{columns}
    % \begin{column}[T]{0.5\linewidth}
      % \begin{itemize}
        % \item \textbf{Push-ups}: Strengthens \textbf{upper body}.
        % \item \textbf{Squats}: Builds \textbf{leg muscles}.
        % \item \textbf{Planks}: Engages \textbf{core muscles}.
        % \item \textbf{Lunges}: Improves \textbf{lower body strength}.
        % \item \textbf{Leg raises}: Strengthens \textbf{abdominal muscles}.
      % \end{itemize}
    % \end{column}
    % \begin{column}[T]{0.5\linewidth}
        % \begin{center}
        % \includegraphics[width=\linewidth,keepaspectratio]{logo_yhk}
        % \end{center}	
    % \end{column}
% \end{columns}
% \end{frame}

%%%%%%%%%%%%%%%%%%%%%%%%%%%%%%%%%%%%%%%%%%%%%%%%%%%%%%%%%%%
\begin{frame}[fragile,shrink=6]\frametitle{Sarvanga Pushti}
\begin{columns}
    \begin{column}[T]{0.5\linewidth}
      \begin{itemize}
		\item \textbf{Sarvanga Pushti}: whole-body practice, done very slowly (Brahmachari)
		\item \textbf{Start}: feet wide apart, fists with thumbs tucked in, one wrist on the other, bend towards one ankle
		\item \textbf{Move}: swing the joined arms up, back and round in a circle with the trunk, then down towards the other ankle
		\item \textbf{Focus}: slow, smooth movement, no jerks; the breath cue is not settled in the sources
		\item \textbf{Benefits}: traditionally credited with strength
		\item \textbf{Cautions}: take expert advice with high blood pressure, heart conditions, back pain or neck pain (a general precaution)
	  \end{itemize}
    \end{column}
    \begin{column}[T]{0.5\linewidth}
		\begin{center}
		% FIG ycb1-fig-sthula-sarvanga | now: images/ycb1_sthula_sarvanga.jpg | final: books/YCB1/images/new/ycb1_sthula_sarvanga_new.png | register: books/YCB1/docs/image_register.csv
		% PROMPT SP-FIG (common text: books/YCB1/docs/image_system_prompt.md) | Used in block ycb1-pd-sthula-vyayama, frame 'Sarvanga Pushti'. Posture: Sarvanga Pushti. One adult practitioner (CS1) in the final held position of Sarvanga Pushti, side view. The whole body in frame with a small margin; head, spine, arms, hands, legs and feet positions exactly as in the notes block for this posture; weight and support points clear. No props unless the notes block names one. No labels. Aspect ratio 4:5 portrait. Check before use: certified teacher confirms the form against the steps in the notes block.
		% NEW-PIC (uncomment to show the new picture beside the original; comment the original only after review) \includegraphics[width=0.35\linewidth,keepaspectratio]{new/ycb1_sthula_sarvanga_new.png}
		\includegraphics[width=\linewidth,keepaspectratio]{ycb1_sthula_sarvanga}
		
		{\tiny (Ref: Day 04 of 30 Days of Yogic Journey - Guiding Principles for Yoga Practitioners and Yogic Sukshma Vyayama - Saatvik Life)}	
		\end{center}	
    \end{column}
\end{columns}
\end{frame}

%%%%%%%%%%%%%%%%%%%%%%%%%%%%%%%%%%%%%%%%%%%%%%%%%%%%%%%%%%%
\begin{frame}[fragile]\frametitle{Hrid Gati (Engine Daud)}\footnotesize
\begin{columns}
    \begin{column}[T]{0.5\linewidth}
      \begin{itemize}
		\item \textbf{Hrid Gati}: running-like practice, the locomotive exercise (Brahmachari)
		\item \textbf{Start}: feet together, elbows bent, fists with thumbs tucked in
		\item \textbf{Move}: run small steps on the toes, heels flicking towards the buttocks, arms moving alternately like pistons; about fifty steps forward, then the same backward
		\item \textbf{Breath}: sharp, deep breaths through the nose, like an engine
		\item \textbf{Benefits}: traditionally credited with endurance and stamina
		\item \textbf{Cautions}: take advice from a qualified teacher or doctor with high blood pressure, heart conditions, back or neck pain (a general precaution)
	  \end{itemize}
    \end{column}
    \begin{column}[T]{0.5\linewidth}
		\begin{center}
		% FIG ycb1-fig-sthula-hrid | now: images/ycb1_sthula_hrid.png | final: books/YCB1/images/new/ycb1_sthula_hrid_new.png | register: books/YCB1/docs/image_register.csv
		% PROMPT SP-FIG (common text: books/YCB1/docs/image_system_prompt.md) | Used in block ycb1-pd-sthula-vyayama, frame 'Hrid Gati (Engine Daud)'. Posture: Hrid Gati (Engine Daud). One adult practitioner (CS1) in the final held position of Hrid Gati (Engine Daud), side view. The whole body in frame with a small margin; head, spine, arms, hands, legs and feet positions exactly as in the notes block for this posture; weight and support points clear. No props unless the notes block names one. No labels. Aspect ratio 4:5 portrait. Check before use: certified teacher confirms the form against the steps in the notes block.
		% NEW-PIC (uncomment to show the new picture beside the original; comment the original only after review) \includegraphics[width=0.35\linewidth,keepaspectratio]{new/ycb1_sthula_hrid_new.png}
		\includegraphics[width=0.5\linewidth,keepaspectratio]{ycb1_sthula_hrid}
		
		{\tiny (Ref: Day 04 of 30 Days of Yogic Journey - Guiding Principles for Yoga Practitioners and Yogic Sukshma Vyayama - Saatvik Life)}	
		\end{center}	
    \end{column}
\end{columns}
\end{frame}

%%%%%%%%%%%%%%%%%%%%%%%%%%%%%%%%%%%%%%%%%%%%%%%%%%%%%%%%%%%
\begin{frame}[fragile]\frametitle{Comparison of Sukshma and Sthula Vyayama}

      \begin{itemize}
        \item \textbf{Sukshma Vyayama}: Focus on \textbf{joints} and \textbf{flexibility}.
        \item \textbf{Sthula Vyayama}: Larger movements of the \textbf{gross body}.
        \item \textbf{Sukshma}: \textbf{Gentle} and \textbf{subtle movements}.
        \item \textbf{Sthula}: \textbf{Dynamic} whole-body movements.
        \item Both \textbf{complement each other} in a balanced Yoga practice.
      \end{itemize}

\end{frame}


%%%%%%%%%%%%%%%%%%%%%%%%%%%%%%%%%%%%%%%%%%%%%%%%%%%%%%%%%%%%%%%%%%%%%%%%%%%%%%%%%%
\begin{frame}[fragile]\frametitle{}
\begin{center}
{\Large Yogic Surya Namaskara}
\end{center}
\end{frame}

%%%%%%%%%%%%%%%%%%%%%%%%%%%%%%%%%%%%%%%%%%%%%%%%%%%%%%%%%%%
\begin{frame}[fragile]\frametitle{Surya Namaskara: Concept}
      \begin{itemize}
        \item \textbf{Surya Namaskara}: salutation to the \textbf{sun}.
        \item Counted as part of \textbf{Hatha Yoga} practice.
        \item A sequence of \textbf{12 postures} joined with the breath; not a single asana.
        \item Loosens the spine and limbs and prepares the body for asanas.
        \item Traditionally done in the early morning or evening facing the sun; custom, not a rule.
		\item Benefits of Surya Namaskara (traditionally credited, for well-being, not cure)
		      \begin{itemize}
			\item \textbf{Flexibility} and \textbf{strength}.
			\item Better \textbf{circulation} and \textbf{digestion}.
			\item \textbf{Calm} and \textbf{focus}.
			\item Help with \textbf{weight control}.
			\item Prepares the body for asana practice.
		\item Precaution and Tips
		  \begin{itemize}
			\item Practise on an \textbf{empty stomach}.
			\item \textbf{Back problems}: only after a yoga expert's advice. Avoid in acute high blood pressure, hernia, heart disease, after a stroke, in fever or inflammation, and during menstruation.
			\item Practise in a \textbf{well-aired} place.
			\item Move slowly and smoothly, never jerky; keep the \textbf{breath steady}, with no forceful breathing or breath holding.
			\item Start with a few rounds and add slowly.
		  \end{itemize}		
		  \end{itemize}
      \end{itemize}

\end{frame}

%%%%%%%%%%%%%%%%%%%%%%%%%%%%%%%%%%%%%%%%%%%%%%%%%%%%%%%%%%%
\begin{frame}[fragile]\frametitle{Sequence of Surya Namaskara}

      \begin{itemize}
        \item \textbf{1. Pranamasana (प्रणामासना)}: Prayer Pose. Breathe normally.
        \item \textbf{2. Hasta Uttanasana (हस्त उत्तानासन)}: Raised Arms Pose (also hastottanasana). Inhale.
        \item \textbf{3. Padahastasana (पदहस्तासन)}: Hand to Foot Pose. Exhale.
        \item \textbf{4. Ashwa Sanchalanasana (अश्व संचारणासन)}: Equestrian Pose. Inhale.
        \item \textbf{5. Parvatasana (पर्वतासन)}: Inverted V Pose (adho mukha svanasana is the common name). Exhale.
        \item \textbf{6. Ashtanga Namaskara (अष्टांग नमस्कार)}: Salute with Eight Points. Exhale; no breath holding.
        \item \textbf{7. Bhujangasana (भुजंगासन)}: Cobra Pose. Inhale.
        \item \textbf{8. Parvatasana (पर्वतासन)}: Inverted V Pose (repeated). Exhale.
        \item \textbf{9. Ashwa Sanchalanasana (अश्व संचारणासन)}: Equestrian Pose (repeated). Inhale.
        \item \textbf{10. Padahastasana (पदहस्तासन)}: Hand to Foot Pose (repeated). Exhale.
        \item \textbf{11. Hasta Uttanasana (हस्त उत्तानासन)}: Raised Arms Pose (repeated). Inhale.
        \item \textbf{12. Pranamasana (प्रणामासना)}: Prayer Pose (repeated). Exhale.
      \end{itemize}
 
\end{frame}


%%%%%%%%%%%%%%%%%%%%%%%%%%%%%%%%%%%%%%%%%%%%%%%%%%%%%%%%%%%
\begin{frame}[fragile]\frametitle{Sequence of Surya Namaskara}

        \begin{center}
        % FIG ycb1-fig-suryanamaskar | now: images/ycb1_suryanamaskar.jpg | final: books/YCB1/images/new/ycb1_suryanamaskar_new.png | register: books/YCB1/docs/image_register.csv
        % PROMPT SP-FIG (common text: books/YCB1/docs/image_system_prompt.md) | Used in block ycb1-pd-surya-namaskara, frame 'Sequence of Surya Namaskara'. The Surya Namaskara sequence: one practitioner (CS1) in a strip of panels, one panel per position in the order given in the notes block, side view, each panel with a thin arrow to the next and a breath marker (in or out) as the notes block says. Panels equal size, white gaps, no text except the breath marker symbols. Aspect ratio 16:9. Check before use: certified teacher confirms the form against the steps in the notes block.
        % NEW-PIC (uncomment to show the new picture beside the original; comment the original only after review) \includegraphics[width=0.35\linewidth,keepaspectratio]{new/ycb1_suryanamaskar_new.png}
        \includegraphics[width=0.8\linewidth,keepaspectratio]{ycb1_suryanamaskar}
        \end{center}	

\end{frame}

%%%%%%%%%%%%%%%%%%%%%%%%%%%%%%%%%%%%%%%%%%%%%%%%%%%%%%%%%%%
\begin{frame}[fragile]\frametitle{Sequence of Surya Namaskara}

        \begin{center}
        % FIG ycb1-fig-sunsal1 | now: images/ycb1_sunsal1.png | final: no image file (see register) | register: books/YCB1/docs/image_register.csv
        % PROMPT SP-NONE (common text: books/YCB1/docs/image_system_prompt.md) | Do not generate a picture. This is a screenshot of a table. Retype it as a LaTeX table (deck) and a Markdown table (notes) from the checked notes block ycb1-pd-surya-namaskara, in own words where it is a list of facts. Used in block ycb1-pd-surya-namaskara, frame 'Sequence of Surya Namaskara'.
        % Screenshot retired; retyped from the notes block (positions 1 to 6).
        % \includegraphics[width=0.5\linewidth,keepaspectratio]{ycb1_sunsal1}
        {\small
        \begin{tabular}{@{}rlll@{}}
        \toprule
        & \textbf{Position} & \textbf{Mantra} & \textbf{Breath} \\
        \midrule
        1 & Pranamasana & Om Mitraya Namah & normal \\
        2 & Hastottanasana & Om Ravaye Namah & in \\
        3 & Padahastasana & Om Suryaya Namah & out \\
        4 & Ashwa sanchalanasana & Om Bhanave Namah & in \\
        5 & Parvatasana & Om Khagaya Namah & out \\
        6 & Ashtanga namaskara & Om Pushne Namah & out \\
        \bottomrule
        \end{tabular}}
        \end{center}


\end{frame}

%%%%%%%%%%%%%%%%%%%%%%%%%%%%%%%%%%%%%%%%%%%%%%%%%%%%%%%%%%%
\begin{frame}[fragile]\frametitle{Sequence of Surya Namaskara}

        \begin{center}
        % FIG ycb1-fig-sunsal2 | now: images/ycb1_sunsal2.png | final: no image file (see register) | register: books/YCB1/docs/image_register.csv
        % PROMPT SP-NONE (common text: books/YCB1/docs/image_system_prompt.md) | Do not generate a picture. This is a screenshot of a table. Retype it as a LaTeX table (deck) and a Markdown table (notes) from the checked notes block ycb1-pd-surya-namaskara, in own words where it is a list of facts. Used in block ycb1-pd-surya-namaskara, frame 'Sequence of Surya Namaskara'.
        % Screenshot retired; retyped from the notes block (positions 7 to 12).
        % \includegraphics[width=0.5\linewidth,keepaspectratio]{ycb1_sunsal2}
        {\small
        \begin{tabular}{@{}rlll@{}}
        \toprule
        & \textbf{Position} & \textbf{Mantra} & \textbf{Breath} \\
        \midrule
        7 & Bhujangasana & Om Hiranyagarbhaya Namah & in \\
        8 & Parvatasana & Om Marichaye Namah & out \\
        9 & Ashwa sanchalanasana & Om Adityaya Namah & in \\
        10 & Padahastasana & Om Savitre Namah & out \\
        11 & Hastottanasana & Om Arkaya Namah & in \\
        12 & Pranamasana & Om Bhaskaraya Namah & out \\
        \bottomrule
        \end{tabular}}

		\vspace{0.5em}
		{\small Twelve names, each said silently with its posture: Mitraya, Ravaye, Suryaya, Bhanave, Khagaya, Pushne, Hiranyagarbhaya, Marichaye, Adityaya, Savitre, Arkaya, Bhaskaraya (each with Namah). Mantras are optional.}

        \end{center}	

\end{frame}

%%%%%%%%%%%%%%%%%%%%%%%%%%%%%%%%%%%%%%%%%%%%%%%%%%%%%%%%%%%%%%%%%%%%%%%%%%%%%%%%%%
\begin{frame}[fragile]\frametitle{}
\begin{center}
{\Large Yogasana}
\end{center}
\end{frame}

%%%%%%%%%%%%%%%%%%%%%%%%%%%%%%%%%%%%%%%%%%%%%%%%%%%%%%%%%%%%%%%%%%%%%%%%%%%%%%%%%%
\begin{frame}[fragile]\frametitle{Tadasana}
\begin{columns}
    \begin{column}[T]{0.5\linewidth}
      \begin{itemize}
        \item Stand with feet together, arms by sides.
        \item Distribute weight evenly on both feet.
        \item Engage thighs and lift chest.
        \item Extend arms overhead, palms facing each other.
        \item Hold the pose and breathe deeply.
        \item \textbf{Benefits:} Improves posture, strengthens legs, and enhances concentration.
        \item \textbf{Contraindications:} Avoid rising on the toes with acute cardiac problems, varicose veins or vertigo (Common Yoga Protocol).
      \end{itemize}
    \end{column}
    \begin{column}[T]{0.5\linewidth}
        \begin{center}
		        % FIG ycb1-fig-tadasan | now: images/ycb1_tadasan.png | final: books/YCB1/images/new/ycb1_tadasan_new.png | register: books/YCB1/docs/image_register.csv
		        % PROMPT SP-FIG (common text: books/YCB1/docs/image_system_prompt.md) | Used in block ycb1-pd-yogasana-practice, frame 'Tadasana'. Posture: Tadasana. One adult practitioner (CS1) in the final held position of Tadasana, side view, standing, facing right. The whole body in frame with a small margin; head, spine, arms, hands, legs and feet positions exactly as in the notes block for this posture; weight and support points clear. No props unless the notes block names one. No labels. Aspect ratio 4:5 portrait. Check before use: certified teacher confirms the form against the steps in the notes block.
		        % NEW-PIC (uncomment to show the new picture beside the original; comment the original only after review) \includegraphics[width=0.35\linewidth,keepaspectratio]{new/ycb1_tadasan_new.png}
		        \includegraphics[width=0.35\linewidth,keepaspectratio]{ycb1_tadasan}
				
				{\tiny (Ref: Certification  of Yoga Professionals Official Guidebook For Level I (Instructor))}		
        \end{center}    
    \end{column}
  \end{columns}
\end{frame}

%%%%%%%%%%%%%%%%%%%%%%%%%%%%%%%%%%%%%%%%%%%%%%%%%%%%%%%%%%%%%%%%%%%%%%%%%%%%%%%%%%
\begin{frame}[fragile]\frametitle{Vrikshasana}
\begin{columns}
    \begin{column}[T]{0.5\linewidth}
      \begin{itemize}
        \item Stand in Tadasana position.
        \item Shift weight to one foot, bend the other knee.
        \item Place the sole of the bent foot on the inner thigh of the standing leg.
        \item Join hands in front of the chest or extend overhead.
        \item Hold the position, focus on balance.
        \item \textbf{Benefits:} Enhances balance, strengthens legs, and improves concentration.
        \item \textbf{Contraindications:} Avoid with arthritis, vertigo or obesity (Common Yoga Protocol); also with knee or ankle injury.
      \end{itemize}
    \end{column}
    \begin{column}[T]{0.5\linewidth}
        \begin{center}
		        % FIG ycb1-fig-vrikshasan | now: images/ycb1_vrikshasan.png | final: books/YCB1/images/new/ycb1_vrikshasan_new.png | register: books/YCB1/docs/image_register.csv
		        % PROMPT SP-FIG (common text: books/YCB1/docs/image_system_prompt.md) | Used in block ycb1-pd-yogasana-practice, frame 'Vrikshasana'. Posture: Vrikshasana. One adult practitioner (CS1) in the final held position of Vrikshasana, side view, standing, facing right. The whole body in frame with a small margin; head, spine, arms, hands, legs and feet positions exactly as in the notes block for this posture; weight and support points clear. No props unless the notes block names one. No labels. Aspect ratio 4:5 portrait. Check before use: certified teacher confirms the form against the steps in the notes block.
		        % NEW-PIC (uncomment to show the new picture beside the original; comment the original only after review) \includegraphics[width=0.35\linewidth,keepaspectratio]{new/ycb1_vrikshasan_new.png}
		        \includegraphics[width=0.4\linewidth,keepaspectratio]{ycb1_vrikshasan}
				
				{\tiny (Ref: Certification  of Yoga Professionals Official Guidebook For Level I (Instructor))}	        
		\end{center}    
    \end{column}
  \end{columns}
\end{frame}

%%%%%%%%%%%%%%%%%%%%%%%%%%%%%%%%%%%%%%%%%%%%%%%%%%%%%%%%%%%%%%%%%%%%%%%%%%%%%%%%%%
\begin{frame}[fragile]\frametitle{Ardha Chakrasana}
\begin{columns}
    \begin{column}[T]{0.5\linewidth}
      \begin{itemize}
        \item Stand with feet shoulder-width apart.
        \item Place hands on lower back for support.
        \item Drop the head back; inhale and bend back from the lower back.
        \item Exhale and relax, then come up slowly with an inhalation.
        \item Hold for 10 to 30 seconds with normal breathing.
        \item \textbf{Benefits:} Makes the spine flexible and strengthens the spinal nerves and neck muscles.
        \item \textbf{Contraindications:} Avoid with vertigo or giddiness; those with high blood pressure bend with care (Common Yoga Protocol).
      \end{itemize}
    \end{column}
    \begin{column}[T]{0.5\linewidth}
        \begin{center}
        \begin{center}
		        % FIG ycb1-fig-ardhachakrasan | now: images/ycb1_ardhachakrasan.png | final: books/YCB1/images/new/ycb1_ardhachakrasan_new.png | register: books/YCB1/docs/image_register.csv
		        % PROMPT SP-FIG (common text: books/YCB1/docs/image_system_prompt.md) | Used in block ycb1-pd-yogasana-practice, frame 'Ardha Chakrasana or Hastottanasan'. Posture: Ardha Chakrasana or Hastottanasan. One adult practitioner (CS1) in the final held position of Ardha Chakrasana or Hastottanasan, side view, standing, facing right. The whole body in frame with a small margin; head, spine, arms, hands, legs and feet positions exactly as in the notes block for this posture; weight and support points clear. No props unless the notes block names one. No labels. Aspect ratio 4:5 portrait. Check before use: certified teacher confirms the form against the steps in the notes block.
		        % NEW-PIC (uncomment to show the new picture beside the original; comment the original only after review) \includegraphics[width=0.35\linewidth,keepaspectratio]{new/ycb1_ardhachakrasan_new.png}
		        \includegraphics[width=\linewidth,keepaspectratio]{ycb1_ardhachakrasan}
				
				{\tiny (Ref: Certification  of Yoga Professionals Official Guidebook For Level I (Instructor))}	        
		\end{center}   
        \end{center}    
    \end{column}
  \end{columns}
\end{frame}

%%%%%%%%%%%%%%%%%%%%%%%%%%%%%%%%%%%%%%%%%%%%%%%%%%%%%%%%%%%%%%%%%%%%%%%%%%%%%%%%%%
\begin{frame}[fragile]\frametitle{Padahastasana}
\begin{columns}
    \begin{column}[T]{0.5\linewidth}
      \begin{itemize}
        \item Stand with feet together, arms by sides.
        \item Inhale and raise arms overhead.
        \item Exhale and bend forward, reaching for the feet.
        \item Keep knees slightly bent if needed.
        \item Hold the pose and breathe deeply.
        \item \textbf{Benefits:} Makes the spine flexible and improves digestion (Common Yoga Protocol).
        \item \textbf{Contraindications:} Avoid with cardiac or back problems, hernia, ulcers, abdominal inflammation, high myopia, vertigo or in pregnancy (Common Yoga Protocol).
      \end{itemize}
    \end{column}
    \begin{column}[T]{0.5\linewidth}
        \begin{center}
        \begin{center}
		        % FIG ycb1-fig-padahastasan | now: images/ycb1_padahastasan.png | final: books/YCB1/images/new/ycb1_padahastasan_new.png | register: books/YCB1/docs/image_register.csv
		        % PROMPT SP-FIG (common text: books/YCB1/docs/image_system_prompt.md) | Used in block ycb1-pd-yogasana-practice, frame 'Padahastasana'. Posture: Padahastasana. One adult practitioner (CS1) in the final held position of Padahastasana, side view, standing, facing right. The whole body in frame with a small margin; head, spine, arms, hands, legs and feet positions exactly as in the notes block for this posture; weight and support points clear. No props unless the notes block names one. No labels. Aspect ratio 4:5 portrait. Check before use: certified teacher confirms the form against the steps in the notes block.
		        % NEW-PIC (uncomment to show the new picture beside the original; comment the original only after review) \includegraphics[width=0.35\linewidth,keepaspectratio]{new/ycb1_padahastasan_new.png}
		        \includegraphics[width=0.9\linewidth,keepaspectratio]{ycb1_padahastasan}
				
				{\tiny (Ref: Certification  of Yoga Professionals Official Guidebook For Level I (Instructor))}	        
		\end{center}   
        \end{center}    
    \end{column}
  \end{columns}
\end{frame}

%%%%%%%%%%%%%%%%%%%%%%%%%%%%%%%%%%%%%%%%%%%%%%%%%%%%%%%%%%%%%%%%%%%%%%%%%%%%%%%%%%
\begin{frame}[fragile,shrink=6]\frametitle{Kati Chakrasana}
\begin{columns}
    \begin{column}[T]{0.5\linewidth}
      \begin{itemize}
        \item Stand with feet about half a metre apart; inhale and raise the arms to shoulder level.
        \item Exhale and twist left: right hand to the left shoulder, left arm round the back to the waist. Look over the left shoulder.
        \item Hold the breath out about 2 seconds, inhale and return; repeat on the other side. Do 5 to 10 rounds, without jerking.
        \item \textbf{Benefits:} Tones the waist, back and hips (Satyananda).
        \item \textbf{Contraindications:} Take care with back or spine problems (general caution; the source states none).
      \end{itemize}
    \end{column}
    \begin{column}[T]{0.5\linewidth}
        \begin{center}
        \begin{center}
		        % FIG ycb1-fig-katichakrasan | now: images/ycb1_katichakrasan.png | final: books/YCB1/images/new/ycb1_katichakrasan_new.png | register: books/YCB1/docs/image_register.csv
		        % PROMPT SP-FIG (common text: books/YCB1/docs/image_system_prompt.md) | Used in block ycb1-pd-yogasana-practice, frame 'Kati Chakrasana'. Posture: Kati Chakrasana. One adult practitioner (CS1) in the final held position of Kati Chakrasana, side view, standing, facing right. The whole body in frame with a small margin; head, spine, arms, hands, legs and feet positions exactly as in the notes block for this posture; weight and support points clear. No props unless the notes block names one. No labels. Aspect ratio 4:5 portrait. Check before use: certified teacher confirms the form against the steps in the notes block.
		        % NEW-PIC (uncomment to show the new picture beside the original; comment the original only after review) \includegraphics[width=0.35\linewidth,keepaspectratio]{new/ycb1_katichakrasan_new.png}
		        \includegraphics[width=0.6\linewidth,keepaspectratio]{ycb1_katichakrasan}
				
				{\tiny (Ref: Prana Yoga)}	        
		\end{center}   
        \end{center}    
    \end{column}
  \end{columns}
\end{frame}

%%%%%%%%%%%%%%%%%%%%%%%%%%%%%%%%%%%%%%%%%%%%%%%%%%%%%%%%%%%%%%%%%%%%%%%%%%%%%%%%%%
\begin{frame}[fragile,shrink=6]\frametitle{Trikonasana}
\begin{columns}
    \begin{column}[T]{0.5\linewidth}
      \begin{itemize}
        \item Stand with feet wide apart, arms extended.
        \item Turn one foot out and the other foot slightly in.
        \item Reach down, placing the hand just behind the foot, or on the ankle, shin or knee if the foot is out of reach.
        \item Extend the other arm upwards, gaze up.
        \item Hold the position, then switch sides.
        \item \textbf{Benefits:} Strengthens calf, thigh and waist muscles and makes the spine flexible.
        \item \textbf{Contraindications:} Avoid with a slipped disc, sciatica or after abdominal surgery; do not overdo the lateral stretch.
      \end{itemize}
    \end{column}
    \begin{column}[T]{0.5\linewidth}
        \begin{center}
        \begin{center}
		        % FIG ycb1-fig-trikonasan | now: images/ycb1_trikonasan.png | final: books/YCB1/images/new/ycb1_trikonasan_new.png | register: books/YCB1/docs/image_register.csv
		        % PROMPT SP-FIG (common text: books/YCB1/docs/image_system_prompt.md) | Used in block ycb1-pd-yogasana-practice, frame 'Trikonasana'. Posture: Trikonasana. One adult practitioner (CS1) in the final held position of Trikonasana, side view, standing, facing right. The whole body in frame with a small margin; head, spine, arms, hands, legs and feet positions exactly as in the notes block for this posture; weight and support points clear. No props unless the notes block names one. No labels. Aspect ratio 4:5 portrait. Check before use: certified teacher confirms the form against the steps in the notes block.
		        % NEW-PIC (uncomment to show the new picture beside the original; comment the original only after review) \includegraphics[width=0.35\linewidth,keepaspectratio]{new/ycb1_trikonasan_new.png}
		        \includegraphics[width=\linewidth,keepaspectratio]{ycb1_trikonasan}
				
				{\tiny (Ref: Certification  of Yoga Professionals Official Guidebook For Level I (Instructor))}	        
		\end{center}   
        \end{center}    
    \end{column}
  \end{columns}
\end{frame}

%%%%%%%%%%%%%%%%%%%%%%%%%%%%%%%%%%%%%%%%%%%%%%%%%%%%%%%%%%%%%%%%%%%%%%%%%%%%%%%%%%
\begin{frame}[fragile]\frametitle{Dandasana}
\begin{columns}
    \begin{column}[T]{0.5\linewidth}
      \begin{itemize}
        \item Sit with both legs stretched out and the heels together.
        \item Palms on the floor beside the buttocks.
        \item Keep spine, neck and head upright; close the eyes.
        \item The starting position for the sitting asanas. Relaxed form: legs a little apart, lean slightly back on the hands.
        \item \textbf{Source:} NIOS (secondary source); no benefits or cautions are stated there.
      \end{itemize}
    \end{column}
    \begin{column}[T]{0.5\linewidth}
        \begin{center}
        \begin{center}
		        % FIG ycb1-fig-dandasan | now: images/ycb1_dandasan.png | final: books/YCB1/images/new/ycb1_dandasan_new.png | register: books/YCB1/docs/image_register.csv
		        % PROMPT SP-FIG (common text: books/YCB1/docs/image_system_prompt.md) | Used in block ycb1-pd-yogasana-practice, frame 'Dandasana'. Posture: Dandasana. One adult practitioner (CS1) in the final held position of Dandasana, front three-quarter view, seated on a plain mat. The whole body in frame with a small margin; head, spine, arms, hands, legs and feet positions exactly as in the notes block for this posture; weight and support points clear. No props unless the notes block names one. No labels. Aspect ratio 4:5 portrait. Check before use: certified teacher confirms the form against the steps in the notes block.
		        % NEW-PIC (uncomment to show the new picture beside the original; comment the original only after review) \includegraphics[width=0.35\linewidth,keepaspectratio]{new/ycb1_dandasan_new.png}
		        \includegraphics[width=\linewidth,keepaspectratio]{ycb1_dandasan}
				
				{\tiny (Ref: Prana Yoga)}	        
		\end{center}   
        \end{center}    
    \end{column}
  \end{columns}
\end{frame}

%%%%%%%%%%%%%%%%%%%%%%%%%%%%%%%%%%%%%%%%%%%%%%%%%%%%%%%%%%%%%%%%%%%%%%%%%%%%%%%%%%
\begin{frame}[fragile]\frametitle{Sukhasana}
\begin{columns}
    \begin{column}[T]{0.5\linewidth}
      \begin{itemize}
        \item Sit with legs crossed comfortably.
        \item Place hands on knees or in a mudra.
        \item Keep spine upright and shoulders relaxed.
        \item Close eyes and focus on breath.
        \item Hold the position, breathing deeply.
        \item \textbf{Benefits:} Promotes relaxation, improves flexibility, and calms the mind.
        \item \textbf{Contraindications:} Avoid if you have knee or hip injuries.
      \end{itemize}
    \end{column}
    \begin{column}[T]{0.5\linewidth}
        \begin{center}
        \begin{center}
		        % FIG ycb1-fig-sukhasan | now: images/ycb1_sukhasan.png | final: books/YCB1/images/new/ycb1_sukhasan_new.png | register: books/YCB1/docs/image_register.csv
		        % PROMPT SP-FIG (common text: books/YCB1/docs/image_system_prompt.md) | Used in block ycb1-pd-yogasana-practice, frame 'Sukhasana'. Posture: Sukhasana. One adult practitioner (CS1) in the final held position of Sukhasana, front three-quarter view, seated on a plain mat. The whole body in frame with a small margin; head, spine, arms, hands, legs and feet positions exactly as in the notes block for this posture; weight and support points clear. No props unless the notes block names one. No labels. Aspect ratio 4:5 portrait. Check before use: certified teacher confirms the form against the steps in the notes block.
		        % NEW-PIC (uncomment to show the new picture beside the original; comment the original only after review) \includegraphics[width=0.35\linewidth,keepaspectratio]{new/ycb1_sukhasan_new.png}
		        \includegraphics[width=\linewidth,keepaspectratio]{ycb1_sukhasan}
				
				{\tiny (Ref: Certification  of Yoga Professionals Official Guidebook For Level I (Instructor))}	        
		\end{center}   
        \end{center}    
    \end{column}
  \end{columns}
\end{frame}

%%%%%%%%%%%%%%%%%%%%%%%%%%%%%%%%%%%%%%%%%%%%%%%%%%%%%%%%%%%%%%%%%%%%%%%%%%%%%%%%%%
\begin{frame}[fragile]\frametitle{Padmasana}
\begin{columns}
    \begin{column}[T]{0.5\linewidth}
      \begin{itemize}
        \item Sit with legs extended, then bend one knee.
        \item Place the foot on the opposite thigh.
        \item Repeat with the other leg, placing the foot on the opposite thigh.
        \item Keep spine straight and shoulders relaxed.
        \item Hold the position, focusing on breath.
        \item \textbf{Benefits:} Used for meditation and pranayama; if full lotus is not possible, use half lotus or sukhasana.
        \item \textbf{Contraindications:} Avoid with a knee or hip injury; never force the knees.
      \end{itemize}
    \end{column}
    \begin{column}[T]{0.5\linewidth}
        \begin{center}
        \begin{center}
		        % FIG ycb1-fig-padmasan | now: images/ycb1_padmasan.png | final: books/YCB1/images/new/ycb1_padmasan_new.png | register: books/YCB1/docs/image_register.csv
		        % PROMPT SP-FIG (common text: books/YCB1/docs/image_system_prompt.md) | Used in block ycb1-pd-yogasana-practice, frame 'Padmasana'. Posture: Padmasana. One adult practitioner (CS1) in the final held position of Padmasana, front three-quarter view, seated on a plain mat. The whole body in frame with a small margin; head, spine, arms, hands, legs and feet positions exactly as in the notes block for this posture; weight and support points clear. No props unless the notes block names one. No labels. Aspect ratio 4:5 portrait. Check before use: certified teacher confirms the form against the steps in the notes block.
		        % NEW-PIC (uncomment to show the new picture beside the original; comment the original only after review) \includegraphics[width=0.35\linewidth,keepaspectratio]{new/ycb1_padmasan_new.png}
		        \includegraphics[width=\linewidth,keepaspectratio]{ycb1_padmasan}
				
				{\tiny (Ref: Certification  of Yoga Professionals Official Guidebook For Level I (Instructor))}	        
		\end{center}   
        \end{center}    
    \end{column}
  \end{columns}
\end{frame}

%%%%%%%%%%%%%%%%%%%%%%%%%%%%%%%%%%%%%%%%%%%%%%%%%%%%%%%%%%%%%%%%%%%%%%%%%%%%%%%%%%
\begin{frame}[fragile,shrink=6]\frametitle{Vajrasana}
\begin{columns}
    \begin{column}[T]{0.5\linewidth}
      \begin{itemize}
        \item Kneel on the floor, sit back on heels.
        \item Keep thighs perpendicular to the floor.
        \item Place hands on knees, palms facing down.
        \item Keep spine straight and shoulders relaxed.
        \item Breathe deeply, holding the position.
        \item \textbf{Benefits:} A steady kneeling posture and the starting position for ardha ushtrasana and shashankasana (Common Yoga Protocol).
        \item \textbf{Contraindications:} Avoid with knee or ankle injuries; people with knee osteoarthritis take care or avoid it (Common Yoga Protocol).
      \end{itemize}
    \end{column}
    \begin{column}[T]{0.5\linewidth}
        \begin{center}
        \begin{center}
		        % FIG ycb1-fig-vajrasan | now: images/ycb1_vajrasan.png | final: books/YCB1/images/new/ycb1_vajrasan_new.png | register: books/YCB1/docs/image_register.csv
		        % PROMPT SP-FIG (common text: books/YCB1/docs/image_system_prompt.md) | Used in block ycb1-pd-yogasana-practice, frame 'Vajrasana'. Posture: Vajrasana. One adult practitioner (CS1) in the final held position of Vajrasana, front three-quarter view, seated on a plain mat. The whole body in frame with a small margin; head, spine, arms, hands, legs and feet positions exactly as in the notes block for this posture; weight and support points clear. No props unless the notes block names one. No labels. Aspect ratio 4:5 portrait. Check before use: certified teacher confirms the form against the steps in the notes block.
		        % NEW-PIC (uncomment to show the new picture beside the original; comment the original only after review) \includegraphics[width=0.35\linewidth,keepaspectratio]{new/ycb1_vajrasan_new.png}
		        \includegraphics[width=0.8\linewidth,keepaspectratio]{ycb1_vajrasan}
				
				{\tiny (Ref: Certification  of Yoga Professionals Official Guidebook For Level I (Instructor))}	        
		\end{center}   
        \end{center}    
    \end{column}
  \end{columns}
\end{frame}

%%%%%%%%%%%%%%%%%%%%%%%%%%%%%%%%%%%%%%%%%%%%%%%%%%%%%%%%%%%%%%%%%%%%%%%%%%%%%%%%%%
\begin{frame}[fragile]\frametitle{Bhadrasana}
\begin{columns}
    \begin{column}[T]{0.5\linewidth}
      \begin{itemize}
        \item Sit with legs extended, then bend knees and bring feet together.
        \item Place feet close to the pelvis, holding toes with hands.
        \item Press the knees gently towards the floor; use a soft cushion under them if the thighs do not reach it.
        \item Keep spine erect and shoulders relaxed.
        \item Hold the pose and breathe deeply.
        \item \textbf{Benefits:} Opens hips, improves flexibility, and calms the mind.
        \item \textbf{Contraindications:} Avoid with severe arthritis or sciatica (Common Yoga Protocol); also knee or hip injury.
      \end{itemize}
    \end{column}
    \begin{column}[T]{0.5\linewidth}
        \begin{center}
        \begin{center}
		        % FIG ycb1-fig-bhadrasan | now: images/ycb1_bhadrasan.png | final: books/YCB1/images/new/ycb1_bhadrasan_new.png | register: books/YCB1/docs/image_register.csv
		        % PROMPT SP-FIG (common text: books/YCB1/docs/image_system_prompt.md) | Used in block ycb1-pd-yogasana-practice, frame 'Bhadrasana'. Posture: Bhadrasana. One adult practitioner (CS1) in the final held position of Bhadrasana, front three-quarter view, seated on a plain mat. The whole body in frame with a small margin; head, spine, arms, hands, legs and feet positions exactly as in the notes block for this posture; weight and support points clear. No props unless the notes block names one. No labels. Aspect ratio 4:5 portrait. Check before use: certified teacher confirms the form against the steps in the notes block.
		        % NEW-PIC (uncomment to show the new picture beside the original; comment the original only after review) \includegraphics[width=0.35\linewidth,keepaspectratio]{new/ycb1_bhadrasan_new.png}
		        \includegraphics[width=\linewidth,keepaspectratio]{ycb1_bhadrasan}
				
				{\tiny (Ref: Patanjali Japan Foundation)}	        
		\end{center}   
        \end{center}    
    \end{column}
  \end{columns}
\end{frame}

%%%%%%%%%%%%%%%%%%%%%%%%%%%%%%%%%%%%%%%%%%%%%%%%%%%%%%%%%%%%%%%%%%%%%%%%%%%%%%%%%%
\begin{frame}[fragile,shrink=6]\frametitle{Mandukasana}
\begin{columns}
    \begin{column}[T]{0.5\linewidth}
      \begin{itemize}
        \item Sit in Vajrasana. Make fists, thumbs tucked inside, and press them against the navel.
        \item Exhale and bend forward, gaze ahead; hold the breath out.
        \item Inhale and come up. Repeat three to four times.
        \item \textbf{Benefits:} Flexibility of thighs and legs.
        \item \textbf{Contraindications:} Avoid with abdominal injury, back pain or knee injury.
        \item Source: CBSE Class XI handbook (secondary source).
      \end{itemize}
    \end{column}
    \begin{column}[T]{0.5\linewidth}
        \begin{center}
        \begin{center}
		        % FIG ycb1-fig-mandukasan | now: images/ycb1_mandukasan.png | final: books/YCB1/images/new/ycb1_mandukasan_new.png | register: books/YCB1/docs/image_register.csv
		        % PROMPT SP-FIG (common text: books/YCB1/docs/image_system_prompt.md) | Used in block ycb1-pd-yogasana-practice, frame 'Mandukasana'. Posture: Mandukasana. One adult practitioner (CS1) in the final held position of Mandukasana, front three-quarter view, seated on a plain mat. The whole body in frame with a small margin; head, spine, arms, hands, legs and feet positions exactly as in the notes block for this posture; weight and support points clear. No props unless the notes block names one. No labels. Aspect ratio 4:5 portrait. Check before use: certified teacher confirms the form against the steps in the notes block.
		        % NEW-PIC (uncomment to show the new picture beside the original; comment the original only after review) \includegraphics[width=0.35\linewidth,keepaspectratio]{new/ycb1_mandukasan_new.png}
		        \includegraphics[width=\linewidth,keepaspectratio]{ycb1_mandukasan}
				
				{\tiny (Ref: Atma Bodh)}	        
		\end{center}   
        \end{center}    
    \end{column}
  \end{columns}
\end{frame}

%%%%%%%%%%%%%%%%%%%%%%%%%%%%%%%%%%%%%%%%%%%%%%%%%%%%%%%%%%%%%%%%%%%%%%%%%%%%%%%%%%
\begin{frame}[fragile]\frametitle{Ushtrasana}
\begin{columns}
    \begin{column}[T]{0.5\linewidth}
      \begin{itemize}
        \item Kneel with knees hip-width apart.
        \item Place hands on lower back for support.
        \item Inhale and lift chest, arching back.
        \item Reach for heels with hands, if possible.
        \item Hold the position, breathing deeply.
        \item \textbf{Benefits:} Stretches the entire front body, opens chest, and improves posture.
        \item \textbf{Contraindications:} Avoid with hernia, abdominal injury, arthritis, vertigo or pregnancy (Common Yoga Protocol); also low blood pressure, insomnia or migraine (Guidebook).
      \end{itemize}
    \end{column}
    \begin{column}[T]{0.5\linewidth}
        \begin{center}
        \begin{center}
		        % FIG ycb1-fig-ushtrasan | now: images/ycb1_ushtrasan.png | final: books/YCB1/images/new/ycb1_ushtrasan_new.png | register: books/YCB1/docs/image_register.csv
		        % PROMPT SP-FIG (common text: books/YCB1/docs/image_system_prompt.md) | Used in block ycb1-pd-yogasana-practice, frame 'Ushtrasana'. Posture: Ushtrasana. One adult practitioner (CS1) in the final held position of Ushtrasana, front three-quarter view, seated on a plain mat. The whole body in frame with a small margin; head, spine, arms, hands, legs and feet positions exactly as in the notes block for this posture; weight and support points clear. No props unless the notes block names one. No labels. Aspect ratio 4:5 portrait. Check before use: certified teacher confirms the form against the steps in the notes block.
		        % NEW-PIC (uncomment to show the new picture beside the original; comment the original only after review) \includegraphics[width=0.35\linewidth,keepaspectratio]{new/ycb1_ushtrasan_new.png}
		        \includegraphics[width=\linewidth,keepaspectratio]{ycb1_ushtrasan}
				
				{\tiny (Ref: Certification  of Yoga Professionals Official Guidebook For Level I (Instructor))}	        
		\end{center}   
        \end{center}    
    \end{column}
  \end{columns}
\end{frame}

%%%%%%%%%%%%%%%%%%%%%%%%%%%%%%%%%%%%%%%%%%%%%%%%%%%%%%%%%%%%%%%%%%%%%%%%%%%%%%%%%%
\begin{frame}[fragile]\frametitle{Shashankasana}
\begin{columns}
    \begin{column}[T]{0.5\linewidth}
      \begin{itemize}
        \item Sit in Vajrasana, then spread the knees wide apart with the big toes touching.
        \item Extend arms forward on the floor.
        \item Exhale, bend forward and rest the head on the ground.
        \item Hold the position, breathing deeply.
        \item \textbf{Benefits:} Relieves stress, stretches back and thighs.
        \item \textbf{Contraindications:} Avoid with acute backache; take care with knee osteoarthritis (Common Yoga Protocol).
      \end{itemize}
    \end{column}
    \begin{column}[T]{0.5\linewidth}
        \begin{center}
        \begin{center}
		        % FIG ycb1-fig-shashankasan | now: images/ycb1_shashankasan.png | final: books/YCB1/images/new/ycb1_shashankasan_new.png | register: books/YCB1/docs/image_register.csv
		        % PROMPT SP-FIG (common text: books/YCB1/docs/image_system_prompt.md) | Used in block ycb1-pd-yogasana-practice, frame 'Shashankasana'. Posture: Shashankasana. One adult practitioner (CS1) in the final held position of Shashankasana, front three-quarter view, seated on a plain mat. The whole body in frame with a small margin; head, spine, arms, hands, legs and feet positions exactly as in the notes block for this posture; weight and support points clear. No props unless the notes block names one. No labels. Aspect ratio 4:5 portrait. Check before use: certified teacher confirms the form against the steps in the notes block.
		        % NEW-PIC (uncomment to show the new picture beside the original; comment the original only after review) \includegraphics[width=0.35\linewidth,keepaspectratio]{new/ycb1_shashankasan_new.png}
		        \includegraphics[width=\linewidth,keepaspectratio]{ycb1_shashankasan}
				
				{\tiny (Ref: Certification  of Yoga Professionals Official Guidebook For Level I (Instructor))}	        
		\end{center}   
        \end{center}    
    \end{column}
  \end{columns}
\end{frame}

%%%%%%%%%%%%%%%%%%%%%%%%%%%%%%%%%%%%%%%%%%%%%%%%%%%%%%%%%%%%%%%%%%%%%%%%%%%%%%%%%%
\begin{frame}[fragile]\frametitle{Uttana Mandukasana}
\begin{columns}
    \begin{column}[T]{0.5\linewidth}
      \begin{itemize}
        \item Upright frog: a sitting posture in Vajrasana with the knees apart and the arms folded behind the head.
        \item Steps: source needed (to be confirmed by a certified teacher).
        \item Not the forward bend with the forehead on the floor.
      \end{itemize}
    \end{column}
    \begin{column}[T]{0.5\linewidth}
        \begin{center}
        \begin{center}
		        % FIG ycb1-fig-uttanamandukasan | now: images/ycb1_uttanamandukasan.jpg | final: books/YCB1/images/new/ycb1_uttanamandukasan_new.png | register: books/YCB1/docs/image_register.csv
		        % PROMPT SP-FIG (common text: books/YCB1/docs/image_system_prompt.md) | Used in block ycb1-pd-yogasana-practice, frame 'Uttana Mandukasana'. Posture: Uttana Mandukasana. One adult practitioner (CS1) in the final held position of Uttana Mandukasana, front three-quarter view, seated on a plain mat. The whole body in frame with a small margin; head, spine, arms, hands, legs and feet positions exactly as in the notes block for this posture; weight and support points clear. No props unless the notes block names one. No labels. Aspect ratio 4:5 portrait. Check before use: certified teacher confirms the form against the steps in the notes block.
		        % NEW-PIC (uncomment to show the new picture beside the original; comment the original only after review) \includegraphics[width=0.35\linewidth,keepaspectratio]{new/ycb1_uttanamandukasan_new.png}
		        \includegraphics[width=0.9\linewidth,keepaspectratio]{ycb1_uttanamandukasan}
				
				{\tiny (Ref: Certification  of Yoga Professionals Official Guidebook For Level I (Instructor))}	        
		\end{center}   
        \end{center}    
    \end{column}
  \end{columns}
\end{frame}

%%%%%%%%%%%%%%%%%%%%%%%%%%%%%%%%%%%%%%%%%%%%%%%%%%%%%%%%%%%%%%%%%%%%%%%%%%%%%%%%%%
\begin{frame}[fragile]\frametitle{Paschimottanasana}
\begin{columns}
    \begin{column}[T]{0.5\linewidth}
      \begin{itemize}
        \item Sit with legs extended, feet flexed.
        \item Inhale and lengthen spine.
        \item Exhale and bend forward, reaching for feet.
        \item Hold the pose and breathe deeply.
        \item \textbf{Benefits:} Stretches the spine and hamstrings, calms the mind.
        \item \textbf{Contraindications:} Avoid if you have back or hamstring injuries.
      \end{itemize}
    \end{column}
    \begin{column}[T]{0.5\linewidth}
        \begin{center}
        \begin{center}
		        % FIG ycb1-fig-pashchimottansasan | now: images/ycb1_pashchimottansasan.png | final: books/YCB1/images/new/ycb1_pashchimottansasan_new.png | register: books/YCB1/docs/image_register.csv
		        % PROMPT SP-FIG (common text: books/YCB1/docs/image_system_prompt.md) | Used in block ycb1-pd-yogasana-practice, frame 'Paschimottanasana'. Posture: Paschimottanasana. One adult practitioner (CS1) in the final held position of Paschimottanasana, front three-quarter view, seated on a plain mat. The whole body in frame with a small margin; head, spine, arms, hands, legs and feet positions exactly as in the notes block for this posture; weight and support points clear. No props unless the notes block names one. No labels. Aspect ratio 4:5 portrait. Check before use: certified teacher confirms the form against the steps in the notes block.
		        % NEW-PIC (uncomment to show the new picture beside the original; comment the original only after review) \includegraphics[width=0.35\linewidth,keepaspectratio]{new/ycb1_pashchimottansasan_new.png}
		        \includegraphics[width=\linewidth,keepaspectratio]{ycb1_pashchimottansasan}
				
				{\tiny (Ref: Certification  of Yoga Professionals Official Guidebook For Level I (Instructor))}	        
		\end{center}   
        \end{center}    
    \end{column}
  \end{columns}
\end{frame}

%%%%%%%%%%%%%%%%%%%%%%%%%%%%%%%%%%%%%%%%%%%%%%%%%%%%%%%%%%%%%%%%%%%%%%%%%%%%%%%%%%
\begin{frame}[fragile,shrink=6]\frametitle{Purvottanasana}
\begin{columns}
    \begin{column}[T]{0.5\linewidth}
      \begin{itemize}
        \item Sit with legs out, feet together; palms on the floor beside the hips, arms straight.
        \item Lean back and lift the hips into one straight line, feet on the floor, head tilted back.
        \item Breathe in lifting, out lowering; steady breath in the hold.
        \item Sources disagree on the direction of the fingers (forward or back); not settled here.
        \item \textbf{Benefits:} Strengthens wrists, arms, shoulders and back; stretches legs and hips.
        \item \textbf{Contraindications:} Avoid with wrist or neck injury.
        \item Source: Art of Living (secondary source).
      \end{itemize}
    \end{column}
    \begin{column}[T]{0.5\linewidth}
        \begin{center}
        \begin{center}
		        % FIG ycb1-fig-purvottansasan | now: images/ycb1_purvottansasan.png | final: books/YCB1/images/new/ycb1_purvottansasan_new.png | register: books/YCB1/docs/image_register.csv
		        % PROMPT SP-FIG (common text: books/YCB1/docs/image_system_prompt.md) | Used in block ycb1-pd-yogasana-practice, frame 'Purvottanasana'. Posture: Purvottanasana. One adult practitioner (CS1) in the final held position of Purvottanasana, front three-quarter view, seated on a plain mat. The whole body in frame with a small margin; head, spine, arms, hands, legs and feet positions exactly as in the notes block for this posture; weight and support points clear. No props unless the notes block names one. No labels. Aspect ratio 4:5 portrait. Check before use: certified teacher confirms the form against the steps in the notes block.
		        % NEW-PIC (uncomment to show the new picture beside the original; comment the original only after review) \includegraphics[width=0.35\linewidth,keepaspectratio]{new/ycb1_purvottansasan_new.png}
		        \includegraphics[width=\linewidth,keepaspectratio]{ycb1_purvottansasan}
				
				{\tiny (Ref: Ekhart Yoga)}	        
		\end{center}   
        \end{center}    
    \end{column}
  \end{columns}
\end{frame}

%%%%%%%%%%%%%%%%%%%%%%%%%%%%%%%%%%%%%%%%%%%%%%%%%%%%%%%%%%%%%%%%%%%%%%%%%%%%%%%%%%
\begin{frame}[fragile,shrink=6]\frametitle{Vakrasana}
\begin{columns}
    \begin{column}[T]{0.5\linewidth}
      \begin{itemize}
        \item Sit with legs extended and back straight.
        \item Bend one leg and place its foot beside the opposite knee.
        \item Exhale and twist towards the bent knee, bringing the opposite arm around the knee to hold the big toe or place the palm beside the foot; the other palm on the floor behind.
        \item Hold for 10 to 30 seconds with normal breathing, then switch sides.
        \item \textbf{Benefits:} Increases spinal flexibility; the Common Yoga Protocol also mentions help with constipation and dyspepsia.
        \item \textbf{Contraindications:} Avoid with severe back pain, vertebral and disc disorders, after abdominal surgery and during menstruation.
      \end{itemize}
    \end{column}
    \begin{column}[T]{0.5\linewidth}
        \begin{center}
        \begin{center}
		        % FIG ycb1-fig-vakrasan | now: images/ycb1_vakrasan.jpg | final: books/YCB1/images/new/ycb1_vakrasan_new.png | register: books/YCB1/docs/image_register.csv
		        % PROMPT SP-FIG (common text: books/YCB1/docs/image_system_prompt.md) | Used in block ycb1-pd-yogasana-practice, frame 'Vakrasana'. Posture: Vakrasana. One adult practitioner (CS1) in the final held position of Vakrasana, front three-quarter view, seated on a plain mat. The whole body in frame with a small margin; head, spine, arms, hands, legs and feet positions exactly as in the notes block for this posture; weight and support points clear. No props unless the notes block names one. No labels. Aspect ratio 4:5 portrait. Check before use: certified teacher confirms the form against the steps in the notes block.
		        % NEW-PIC (uncomment to show the new picture beside the original; comment the original only after review) \includegraphics[width=0.35\linewidth,keepaspectratio]{new/ycb1_vakrasan_new.png}
		        \includegraphics[width=\linewidth,keepaspectratio]{ycb1_vakrasan}
				
				{\tiny (Ref: CONDE NAST TRAVELLER)}	        
		\end{center}   
        \end{center}    
    \end{column}
  \end{columns}
\end{frame}

%%%%%%%%%%%%%%%%%%%%%%%%%%%%%%%%%%%%%%%%%%%%%%%%%%%%%%%%%%%%%%%%%%%%%%%%%%%%%%%%%%
\begin{frame}[fragile,shrink=6]\frametitle{Gomukhasana}
\begin{columns}
    \begin{column}[T]{0.5\linewidth}
      \begin{itemize}
        \item Sit with the right knee directly above the left.
        \item Left arm behind the back, right arm over the shoulder; try to clasp the fingers.
        \item Raised elbow behind the head, spine upright, breathe normally.
        \item Hold up to 2 minutes (Satyananda), 1 minute (NIOS); then swap sides.
        \item \textbf{Benefits:} Loosens shoulders and neck, opens the chest, improves posture.
        \item \textbf{Contraindications:} Avoid with shoulder or knee injury; with bleeding piles.
      \end{itemize}
    \end{column}
    \begin{column}[T]{0.5\linewidth}
        \begin{center}
        \begin{center}
		        % FIG ycb1-fig-gomukhasan | now: images/ycb1_gomukhasan.jpg | final: books/YCB1/images/new/ycb1_gomukhasan_new.png | register: books/YCB1/docs/image_register.csv
		        % PROMPT SP-FIG (common text: books/YCB1/docs/image_system_prompt.md) | Used in block ycb1-pd-yogasana-practice, frame 'Gomukhasana'. Posture: Gomukhasana. One adult practitioner (CS1) in the final held position of Gomukhasana, front three-quarter view, seated on a plain mat. The whole body in frame with a small margin; head, spine, arms, hands, legs and feet positions exactly as in the notes block for this posture; weight and support points clear. No props unless the notes block names one. No labels. Aspect ratio 4:5 portrait. Check before use: certified teacher confirms the form against the steps in the notes block.
		        % NEW-PIC (uncomment to show the new picture beside the original; comment the original only after review) \includegraphics[width=0.35\linewidth,keepaspectratio]{new/ycb1_gomukhasan_new.png}
		        \includegraphics[width=\linewidth,keepaspectratio]{ycb1_gomukhasan}
				
				{\tiny (Ref: Himalayan Yoga Ashram)}	        
		\end{center}   
        \end{center}    
    \end{column}
  \end{columns}
\end{frame}

%%%%%%%%%%%%%%%%%%%%%%%%%%%%%%%%%%%%%%%%%%%%%%%%%%%%%%%%%%%%%%%%%%%%%%%%%%%%%%%%%%
\begin{frame}[fragile]\frametitle{Bhujangasana}
\begin{columns}
    \begin{column}[T]{0.5\linewidth}
      \begin{itemize}
        \item Lie on your stomach, legs extended, and feet together.
        \item Place hands under shoulders, elbows close to the body.
        \item Inhale and slowly lift the chin and chest up to the navel region.
        \item Hold the pose and breathe deeply.
        \item \textbf{Benefits:} Helps with stress management and constipation; strengthens the back (Common Yoga Protocol).
        \item \textbf{Contraindications:} Avoid for two to three months after abdominal surgery, with hernia or ulcers, or with back or wrist injury.
      \end{itemize}
    \end{column}
    \begin{column}[T]{0.5\linewidth}
        \begin{center}
        \begin{center}
		        % FIG ycb1-fig-bhujangasan | now: images/ycb1_bhujangasan.png | final: books/YCB1/images/new/ycb1_bhujangasan_new.png | register: books/YCB1/docs/image_register.csv
		        % PROMPT SP-FIG (common text: books/YCB1/docs/image_system_prompt.md) | Used in block ycb1-pd-yogasana-practice, frame 'Bhujangasana'. Posture: Bhujangasana. One adult practitioner (CS1) in the final held position of Bhujangasana, side view, lying face down on a plain mat. The whole body in frame with a small margin; head, spine, arms, hands, legs and feet positions exactly as in the notes block for this posture; weight and support points clear. No props unless the notes block names one. No labels. Aspect ratio 3:2 landscape. Check before use: certified teacher confirms the form against the steps in the notes block.
		        % NEW-PIC (uncomment to show the new picture beside the original; comment the original only after review) \includegraphics[width=0.35\linewidth,keepaspectratio]{new/ycb1_bhujangasan_new.png}
		        \includegraphics[width=\linewidth,keepaspectratio]{ycb1_bhujangasan}
				
				{\tiny (Ref: Certification  of Yoga Professionals Official Guidebook For Level I (Instructor))}	        
		\end{center}   
        \end{center}    
    \end{column}
  \end{columns}
\end{frame}

%%%%%%%%%%%%%%%%%%%%%%%%%%%%%%%%%%%%%%%%%%%%%%%%%%%%%%%%%%%%%%%%%%%%%%%%%%%%%%%%%%
\begin{frame}[fragile]\frametitle{Shalabhasana}
\begin{columns}
    \begin{column}[T]{0.5\linewidth}
      \begin{itemize}
        \item Lie on your stomach in Makarasana, chin on the floor, hands beside the body with palms up.
        \item Inhale and raise the legs off the floor without bending the knees.
        \item Squeeze the buttocks and pull up the knee caps; stay for 10 to 20 seconds.
        \item Hold the position, breathing deeply.
        \item \textbf{Benefits:} Strengthens lower back, improves posture.
        \item \textbf{Contraindications:} Avoid with cardiac problems, high blood pressure, peptic ulcers or hernia; take care with severe lower back pain.
      \end{itemize}
    \end{column}
    \begin{column}[T]{0.5\linewidth}
        \begin{center}
        \begin{center}
		        % FIG ycb1-fig-shalabhasan | now: images/ycb1_shalabhasan.png | final: books/YCB1/images/new/ycb1_shalabhasan_new.png | register: books/YCB1/docs/image_register.csv
		        % PROMPT SP-FIG (common text: books/YCB1/docs/image_system_prompt.md) | Used in block ycb1-pd-yogasana-practice, frame 'Shalabhasana'. Posture: Shalabhasana. One adult practitioner (CS1) in the final held position of Shalabhasana, side view, lying face down on a plain mat. The whole body in frame with a small margin; head, spine, arms, hands, legs and feet positions exactly as in the notes block for this posture; weight and support points clear. No props unless the notes block names one. No labels. Aspect ratio 3:2 landscape. Check before use: certified teacher confirms the form against the steps in the notes block.
		        % NEW-PIC (uncomment to show the new picture beside the original; comment the original only after review) \includegraphics[width=0.35\linewidth,keepaspectratio]{new/ycb1_shalabhasan_new.png}
		        \includegraphics[width=\linewidth,keepaspectratio]{ycb1_shalabhasan}
				
				{\tiny (Ref: Certification  of Yoga Professionals Official Guidebook For Level I (Instructor))}	        
		\end{center}   
        \end{center}    
    \end{column}
  \end{columns}
\end{frame}

%%%%%%%%%%%%%%%%%%%%%%%%%%%%%%%%%%%%%%%%%%%%%%%%%%%%%%%%%%%%%%%%%%%%%%%%%%%%%%%%%%
\begin{frame}[fragile]\frametitle{Makarasana}
\begin{columns}
    \begin{column}[T]{0.5\linewidth}
      \begin{itemize}
        \item Lie on your stomach with the feet wide apart, toes pointing outward.
        \item Bend both arms and place one hand on the other.
        \item Rest the forehead on the hands and close the eyes.
        \item Breathe deeply and relax.
        \item \textbf{Benefits:} Relaxes the lower back; the relaxation posture for the prone asanas.
        \item \textbf{Contraindications:} Avoid with low blood pressure, severe cardiac problems or pregnancy (Common Yoga Protocol).
      \end{itemize}
    \end{column}
    \begin{column}[T]{0.5\linewidth}
        \begin{center}
        \begin{center}
		        % FIG ycb1-fig-makarasan | now: images/ycb1_makarasan.png | final: books/YCB1/images/new/ycb1_makarasan_new.png | register: books/YCB1/docs/image_register.csv
		        % PROMPT SP-FIG (common text: books/YCB1/docs/image_system_prompt.md) | Used in block ycb1-pd-yogasana-practice, frame 'Makarasana'. Posture: Makarasana. One adult practitioner (CS1) in the final held position of Makarasana, side view, lying face down on a plain mat. The whole body in frame with a small margin; head, spine, arms, hands, legs and feet positions exactly as in the notes block for this posture; weight and support points clear. No props unless the notes block names one. No labels. Aspect ratio 3:2 landscape. Check before use: certified teacher confirms the form against the steps in the notes block.
		        % NEW-PIC (uncomment to show the new picture beside the original; comment the original only after review) \includegraphics[width=0.35\linewidth,keepaspectratio]{new/ycb1_makarasan_new.png}
		        \includegraphics[width=\linewidth,keepaspectratio]{ycb1_makarasan}
				
				{\tiny (Ref: Vydya Health)}	        
		\end{center}   
        \end{center}    
    \end{column}
  \end{columns}
\end{frame}

%%%%%%%%%%%%%%%%%%%%%%%%%%%%%%%%%%%%%%%%%%%%%%%%%%%%%%%%%%%%%%%%%%%%%%%%%%%%%%%%%%
\begin{frame}[fragile]\frametitle{Pavanamuktasana}
\begin{columns}
    \begin{column}[T]{0.5\linewidth}
      \begin{itemize}
        \item Lie on your back, knees bent, and feet on the floor.
        \item Hug knees to chest, interlace fingers around shins.
        \item Exhale and lift the head until the chin or nose nears the knees.
        \item Hold the position and breathe deeply.
        \item \textbf{Benefits:} Helps remove flatulence and constipation and stretches the lumbar region (Common Yoga Protocol).
        \item \textbf{Contraindications:} Avoid with abdominal injuries, hernia, sciatica, severe back pain or in pregnancy.
      \end{itemize}
    \end{column}
    \begin{column}[T]{0.5\linewidth}
        \begin{center}
        \begin{center}
		        % FIG ycb1-fig-pavanamuktasan | now: images/ycb1_pavanamuktasan.png | final: books/YCB1/images/new/ycb1_pavanamuktasan_new.png | register: books/YCB1/docs/image_register.csv
		        % PROMPT SP-FIG (common text: books/YCB1/docs/image_system_prompt.md) | Used in block ycb1-pd-yogasana-practice, frame 'Pavanamuktasana'. Posture: Pavanamuktasana. One adult practitioner (CS1) in the final held position of Pavanamuktasana, side view, lying on the back on a plain mat. The whole body in frame with a small margin; head, spine, arms, hands, legs and feet positions exactly as in the notes block for this posture; weight and support points clear. No props unless the notes block names one. No labels. Aspect ratio 3:2 landscape. Check before use: certified teacher confirms the form against the steps in the notes block.
		        % NEW-PIC (uncomment to show the new picture beside the original; comment the original only after review) \includegraphics[width=0.35\linewidth,keepaspectratio]{new/ycb1_pavanamuktasan_new.png}
		        \includegraphics[width=\linewidth,keepaspectratio]{ycb1_pavanamuktasan}
				
				{\tiny (Ref: Certification  of Yoga Professionals Official Guidebook For Level I (Instructor))}	        
		\end{center}   
        \end{center}    
    \end{column}
  \end{columns}
\end{frame}

%%%%%%%%%%%%%%%%%%%%%%%%%%%%%%%%%%%%%%%%%%%%%%%%%%%%%%%%%%%%%%%%%%%%%%%%%%%%%%%%%%
\begin{frame}[fragile]\frametitle{Uttanapadasana / Ardha Halasana}
\begin{columns}
    \begin{column}[T]{0.5\linewidth}
      \begin{itemize}
        \item Lie on your back, legs extended, arms by the sides. The Board syllabus lists uttanapadasana and ardha halasana as two separate postures; the steps here cover the raised-leg form.
        \item Inhale and raise the legs, knees straight, to a comfortable height.
        \item Keep back and shoulders on the floor.
        \item Hold the position and breathe deeply.
        \item \textbf{Benefits:} Strengthens abdominal muscles, tones legs.
        \item \textbf{Contraindications:} Avoid if you have back or leg issues.
      \end{itemize}
    \end{column}
    \begin{column}[T]{0.5\linewidth}
        \begin{center}
        \begin{center}
		        % FIG ycb1-fig-uttanpadasan | now: images/ycb1_uttanpadasan.jpg | final: books/YCB1/images/new/ycb1_uttanpadasan_new.png | register: books/YCB1/docs/image_register.csv
		        % PROMPT SP-FIG (common text: books/YCB1/docs/image_system_prompt.md) | Used in block ycb1-pd-yogasana-practice, frame 'Uttanapadasana'. Posture: Uttanapadasana. One adult practitioner (CS1) in the final held position of Uttanapadasana, side view, lying on the back on a plain mat. The whole body in frame with a small margin; head, spine, arms, hands, legs and feet positions exactly as in the notes block for this posture; weight and support points clear. No props unless the notes block names one. No labels. Aspect ratio 3:2 landscape. Check before use: certified teacher confirms the form against the steps in the notes block.
		        % NEW-PIC (uncomment to show the new picture beside the original; comment the original only after review) \includegraphics[width=0.35\linewidth,keepaspectratio]{new/ycb1_uttanpadasan_new.png}
		        \includegraphics[width=\linewidth,keepaspectratio]{ycb1_uttanpadasan}
				
				{\tiny (Ref: Bodhi School of Yoga)}	        
		\end{center}   
        \end{center}    
    \end{column}
  \end{columns}
\end{frame}

% %%%%%%%%%%%%%%%%%%%%%%%%%%%%%%%%%%%%%%%%%%%%%%%%%%%%%%%%%%%%%%%%%%%%%%%%%%%%%%%%%%
% \begin{frame}[fragile]\frametitle{}
% \begin{columns}
    % \begin{column}[T]{0.5\linewidth}
      % \begin{itemize}
        % \item Lie on your back, arms by sides.
        % \item Lift legs overhead, bringing toes to the floor behind the head.
        % \item Support lower back with hands if needed.
        % \item Hold the position and breathe deeply.
        % \item \textbf{Benefits:} Stretches back and legs, strengthens core.
        % \item \textbf{Contraindications:} Avoid if you have neck or back issues.
      % \end{itemize}
    % \end{column}
    % \begin{column}[T]{0.5\linewidth}
        % \begin{center}
        % \begin{center}
		        % \includegraphics[width=\linewidth,keepaspectratio]{ycb1_vrikshasan}
				
				% {\tiny (Ref: Certification  of Yoga Professionals Official Guidebook For Level I (Instructor))}	        
		% \end{center}   
        % \end{center}    
    % \end{column}
  % \end{columns}
% \end{frame}

%%%%%%%%%%%%%%%%%%%%%%%%%%%%%%%%%%%%%%%%%%%%%%%%%%%%%%%%%%%%%%%%%%%%%%%%%%%%%%%%%%
\begin{frame}[fragile,shrink=6]\frametitle{Setubandhasana}
\begin{columns}
    \begin{column}[T]{0.5\linewidth}
      \begin{itemize}
        \item Lie on your back, knees bent, heels near the buttocks; hold both ankles.
        \item Lift hips towards the ceiling, pressing into feet.
        \item Keep head and shoulders on the floor; the hands may support the waist (a variant interlaces the fingers under the back).
        \item Hold the position and breathe deeply.
        \item \textbf{Benefits:} Strengthens the lower back and helps digestion and constipation (Common Yoga Protocol).
        \item \textbf{Contraindications:} Avoid with ulcers or hernia and in late pregnancy (Common Yoga Protocol); also neck or back injuries.
      \end{itemize}
    \end{column}
    \begin{column}[T]{0.5\linewidth}
        \begin{center}
        \begin{center}
		        % FIG ycb1-fig-setubandhasan | now: images/ycb1_setubandhasan.jpg | final: books/YCB1/images/new/ycb1_setubandhasan_new.png | register: books/YCB1/docs/image_register.csv
		        % PROMPT SP-FIG (common text: books/YCB1/docs/image_system_prompt.md) | Used in block ycb1-pd-yogasana-practice, frame 'Setubandhasana'. Posture: Setubandhasana. One adult practitioner (CS1) in the final held position of Setubandhasana, side view, lying on the back on a plain mat. The whole body in frame with a small margin; head, spine, arms, hands, legs and feet positions exactly as in the notes block for this posture; weight and support points clear. No props unless the notes block names one. No labels. Aspect ratio 3:2 landscape. Check before use: certified teacher confirms the form against the steps in the notes block.
		        % NEW-PIC (uncomment to show the new picture beside the original; comment the original only after review) \includegraphics[width=0.35\linewidth,keepaspectratio]{new/ycb1_setubandhasan_new.png}
		        \includegraphics[width=\linewidth,keepaspectratio]{ycb1_setubandhasan}
				
				{\tiny (Ref: Rishikesh Yogis Yogashala)}	        
		\end{center}   
        \end{center}    
    \end{column}
  \end{columns}
\end{frame}

%%%%%%%%%%%%%%%%%%%%%%%%%%%%%%%%%%%%%%%%%%%%%%%%%%%%%%%%%%%%%%%%%%%%%%%%%%%%%%%%%%
\begin{frame}[fragile,shrink=6]\frametitle{Vipareetakarani}
\begin{columns}
    \begin{column}[T]{0.5\linewidth}
      \begin{itemize}
        \item The Satyananda hands-on-hips version, not legs-up-the-wall.
        \item Lie flat; raise the legs, lift the hips and cup them with the hands, elbows bent.
        \item Legs vertical, trunk at about 45 degrees; weight on shoulders and elbows, chin off the chest.
        \item Breathe normally. Hold a few seconds at first, building to 3 to 5 minutes.
        \item \textbf{Contraindications:} As for Sarvangasana: neck disc problems, high blood pressure, heart problems, weak eye vessels, enlarged thyroid, liver or spleen; avoid in menstruation and late pregnancy.
      \end{itemize}
    \end{column}
    \begin{column}[T]{0.5\linewidth}
        \begin{center}
        \begin{center}
		        % FIG ycb1-fig-vipartikarni | now: images/ycb1_vipartikarni.png | final: books/YCB1/images/new/ycb1_vipartikarni_new.png | register: books/YCB1/docs/image_register.csv
		        % PROMPT SP-FIG (common text: books/YCB1/docs/image_system_prompt.md) | Used in block ycb1-pd-yogasana-practice, frame 'Vipareetakarani'. Posture: Vipareetakarani. One adult practitioner (CS1) in the final held position of Vipareetakarani, side view, lying on the back on a plain mat. The whole body in frame with a small margin; head, spine, arms, hands, legs and feet positions exactly as in the notes block for this posture; weight and support points clear. No props unless the notes block names one. No labels. Aspect ratio 3:2 landscape. Check before use: certified teacher confirms the form against the steps in the notes block.
		        % NEW-PIC (uncomment to show the new picture beside the original; comment the original only after review) \includegraphics[width=0.35\linewidth,keepaspectratio]{new/ycb1_vipartikarni_new.png}
		        \includegraphics[width=\linewidth,keepaspectratio]{ycb1_vipartikarni}
				
				{\tiny (Ref: Yoga4Lyf)}	        
		\end{center}   
        \end{center}    
    \end{column}
  \end{columns}
\end{frame}

%%%%%%%%%%%%%%%%%%%%%%%%%%%%%%%%%%%%%%%%%%%%%%%%%%%%%%%%%%%%%%%%%%%%%%%%%%%%%%%%%%
\begin{frame}[fragile,shrink=6]\frametitle{Saral Matsyasana}
\begin{columns}
    \begin{column}[T]{0.5\linewidth}
      \begin{itemize}
        \item Easy variation of Matsyasana (Satyananda), legs straight. Name matched by us, to be confirmed by a certified teacher.
        \item Stretch both legs straight and lean back on the arms.
        \item Rest the top of the head on the floor and arch the back; hands on thighs or floor.
        \item Return after a short while.
        \item \textbf{Contraindications:} Avoid with heart disease, ulcer, hernia, back conditions, serious illness or pregnancy; take care with blood pressure or a delicate neck.
      \end{itemize}
    \end{column}
    \begin{column}[T]{0.5\linewidth}
        \begin{center}
        \begin{center}
		        % FIG ycb1-fig-saralmatsyasan | now: images/ycb1_saralmatsyasan.jpg | final: books/YCB1/images/new/ycb1_saralmatsyasan_new.png | register: books/YCB1/docs/image_register.csv
		        % PROMPT SP-FIG (common text: books/YCB1/docs/image_system_prompt.md) | Used in block ycb1-pd-yogasana-practice, frame 'Saral Matsyasana'. Posture: Saral Matsyasana. One adult practitioner (CS1) in the final held position of Saral Matsyasana, side view, lying on the back on a plain mat. The whole body in frame with a small margin; head, spine, arms, hands, legs and feet positions exactly as in the notes block for this posture; weight and support points clear. No props unless the notes block names one. No labels. Aspect ratio 3:2 landscape. Check before use: certified teacher confirms the form against the steps in the notes block.
		        % NEW-PIC (uncomment to show the new picture beside the original; comment the original only after review) \includegraphics[width=0.35\linewidth,keepaspectratio]{new/ycb1_saralmatsyasan_new.png}
		        \includegraphics[width=\linewidth,keepaspectratio]{ycb1_saralmatsyasan}
				
				{\tiny (Ref: Kerala Tourism)}	        
		\end{center}   
        \end{center}    
    \end{column}
  \end{columns}
\end{frame}




%%%%%%%%%%%%%%%%%%%%%%%%%%%%%%%%%%%%%%%%%%%%%%%%%%%%%%%%%%%%%%%%%%%%%%%%%%%%%%%%%%
\begin{frame}[fragile]\frametitle{}
\begin{center}
{\Large Preparatory Breathing Practices }
\end{center}
\end{frame}

%%%%%%%%%%%%%%%%%%%%%%%%%%%%%%%%%%%%%%%%%%%%%%%%%%%%%%%%%%%
\begin{frame}[fragile]\frametitle{Preparatory Breathing Practices: Overview}
      \begin{itemize}
        \item \textbf{Preparatory Breathing}: makes the breath \textbf{slow}, \textbf{deep} and \textbf{smooth} before pranayama.
        \item Sit or lie with head, neck and spine in one line; eyes closed; breathe through the \textbf{nose}.
        \item Inhalation is active; exhalation is passive and relaxing. Never force; stop if dizzy.
        \item Correct diaphragm movement brings better \textbf{oxygen intake} without effort.
        \item The three sections join into the \textbf{full yogic breath}.
      \end{itemize}

\end{frame}

%%%%%%%%%%%%%%%%%%%%%%%%%%%%%%%%%%%%%%%%%%%%%%%%%%%%%%%%%%%
\begin{frame}[fragile]\frametitle{Abdominal Breathing (Diaphragmatic Breathing)}
\begin{columns}
    \begin{column}[T]{0.5\linewidth}
      \begin{itemize}
        \item Palms on the belly; the focus is the \textbf{diaphragm}.
        \item \textbf{Inhale} deeply through the nose; the \textbf{abdomen} swells without pressure.
        \item \textbf{Exhale} slowly through the nose; exhalation is passive and the \textbf{abdomen} gently flattens.
        \item Supports \textbf{relaxation} and eases \textbf{stress}.
        \item The chest stays still; do not draw the abdomen in while inhaling.
      \end{itemize}
    \end{column}
    \begin{column}[T]{0.5\linewidth}
        \begin{center}
        % FIG ycb1-fig-abdominalbreathing | now: images/ycb1_abdominalbreathing.png | final: books/YCB1/images/new/ycb1_abdominalbreathing_new.pdf | register: books/YCB1/docs/image_register.csv
        % PROMPT SP-DIAGRAM (common text: books/YCB1/docs/image_system_prompt.md) | Used in block ycb1-pd-preparatory-breathing, frame 'Abdominal Breathing (Diaphragmatic Breathing)'. Concept diagram: Abdominal Breathing (Diaphragmatic Breathing). Boxes, circles and arrows only, drawn from the structure given in the notes block (the names and their order are listed there); Sanskrit terms in Roman letters with the English meaning in brackets, nothing else added. Flat colours that still read in grayscale. Best redrawn later in TikZ. Aspect ratio 4:3. Check before use: author and the Sanskrit check in the notes.
        % NEW-PIC (uncomment to show the new picture beside the original; comment the original only after review) \includegraphics[width=0.35\linewidth,keepaspectratio]{new/ycb1_abdominalbreathing_new.pdf}
        \includegraphics[width=\linewidth,keepaspectratio]{ycb1_abdominalbreathing}
				
		{\tiny (Ref: Certification  of Yoga Professionals Official Guidebook For Level I (Instructor))}	  		
        \end{center}	
    \end{column}
\end{columns}
\end{frame}

%%%%%%%%%%%%%%%%%%%%%%%%%%%%%%%%%%%%%%%%%%%%%%%%%%%%%%%%%%%
\begin{frame}[fragile]\frametitle{Chest Breathing}
\begin{columns}
    \begin{column}[T]{0.5\linewidth}
      \begin{itemize}
        \item Involves the \textbf{chest} and \textbf{rib cage}; the ribs widen in all directions.
        \item \textbf{Inhale} to expand the \textbf{chest} and \textbf{rib cage}.
        \item \textbf{Exhale} as the ribs return; abdomen still, shoulders relaxed.
        \item Takes in less air than abdominal breathing and becomes dominant under stress.
        \item Unites with the other two sections in the \textbf{full yogic breath}.
      \end{itemize}
    \end{column}
    \begin{column}[T]{0.5\linewidth}
        \begin{center}
        % FIG ycb1-fig-chestbreathing | now: images/ycb1_chestbreathing.png | final: books/YCB1/images/new/ycb1_chestbreathing_new.pdf | register: books/YCB1/docs/image_register.csv
        % PROMPT SP-DIAGRAM (common text: books/YCB1/docs/image_system_prompt.md) | Used in block ycb1-pd-preparatory-breathing, frame 'Chest Breathing'. Concept diagram: Chest Breathing. Boxes, circles and arrows only, drawn from the structure given in the notes block (the names and their order are listed there); Sanskrit terms in Roman letters with the English meaning in brackets, nothing else added. Flat colours that still read in grayscale. Best redrawn later in TikZ. Aspect ratio 4:3. Check before use: author and the Sanskrit check in the notes.
        % NEW-PIC (uncomment to show the new picture beside the original; comment the original only after review) \includegraphics[width=0.35\linewidth,keepaspectratio]{new/ycb1_chestbreathing_new.pdf}
        \includegraphics[width=\linewidth,keepaspectratio]{ycb1_chestbreathing}
				
		{\tiny (Ref: Certification  of Yoga Professionals Official Guidebook For Level I (Instructor))}	  
        \end{center}	
    \end{column}
\end{columns}
\end{frame}

%%%%%%%%%%%%%%%%%%%%%%%%%%%%%%%%%%%%%%%%%%%%%%%%%%%%%%%%%%%
\begin{frame}[fragile]\frametitle{Clavicular Breathing}
\begin{columns}
    \begin{column}[T]{0.5\linewidth}
      \begin{itemize}
        \item Focuses on \textbf{upper chest} and \textbf{collarbone}.
        \item \textbf{Inhale} to lift the \textbf{clavicles} and expand the \textbf{upper chest}.
        \item \textbf{Exhale} to lower the \textbf{clavicles}.
        \item A \textbf{shallow} breath; keep the shoulders still.
        \item The last step before the full yogic breath; practised as a step, not on its own.
      \end{itemize}
    \end{column}
    \begin{column}[T]{0.5\linewidth}
        \begin{center}
        % FIG ycb1-fig-clavicularbreathing | now: images/ycb1_clavicularbreathing.jpg | final: books/YCB1/images/new/ycb1_clavicularbreathing_new.png | register: books/YCB1/docs/image_register.csv
        % PROMPT SP-FIG (common text: books/YCB1/docs/image_system_prompt.md) | Used in block ycb1-pd-preparatory-breathing, frame 'Clavicular Breathing'. Preparatory breathing: Clavicular Breathing. One adult practitioner (CS1), side view, seated or lying as the notes block says, hands placed as the notes block says. Show the torso with a soft translucent highlight on the part that rises on the in-breath and a thin arrow for the direction of movement; a faint outline of the lungs and diaphragm may sit behind the highlight. Calm face, eyes closed. Aspect ratio 4:3 landscape. Check before use: certified teacher confirms the form against the steps in the notes block.
        % NEW-PIC (uncomment to show the new picture beside the original; comment the original only after review) \includegraphics[width=0.35\linewidth,keepaspectratio]{new/ycb1_clavicularbreathing_new.png}
        \includegraphics[width=\linewidth,keepaspectratio]{ycb1_clavicularbreathing}
				
		{\tiny (Ref: Sri Sri School of Yoga)}	  
        \end{center}	
    \end{column}
\end{columns}
\end{frame}

%%%%%%%%%%%%%%%%%%%%%%%%%%%%%%%%%%%%%%%%%%%%%%%%%%%%%%%%%%%
\begin{frame}[fragile]\frametitle{Combination Breathing (Three-Part Breathing)}
\begin{columns}
    \begin{column}[T]{0.5\linewidth}
      \begin{itemize}
        \item Combines \textbf{abdominal}, \textbf{chest}, and \textbf{clavicular} breathing.
        \item \textbf{Inhale} first into the abdomen, then the chest, and finally the clavicles.
        \item \textbf{Exhale} in reverse order, as one smooth wave.
        \item Uses the \textbf{lungs fully}; also called the \textbf{full yogic breath}.
        \item Begin with 5 to 10 rounds; if dizzy, stop and breathe normally.
      \end{itemize}
    \end{column}
    \begin{column}[T]{0.5\linewidth}
        \begin{center}
        % FIG ycb1-fig-threepartbreathing | now: images/ycb1_threepartbreathing.jpg | final: books/YCB1/images/new/ycb1_threepartbreathing_new.png | register: books/YCB1/docs/image_register.csv
        % PROMPT SP-FIG (common text: books/YCB1/docs/image_system_prompt.md) | Used in block ycb1-pd-preparatory-breathing, frame 'Combination Breathing (Three-Part Breathing)'. Preparatory breathing: Combination Breathing (Three-Part Breathing). One adult practitioner (CS1), side view, seated or lying as the notes block says, hands placed as the notes block says. Show the torso with a soft translucent highlight on the part that rises on the in-breath and a thin arrow for the direction of movement; a faint outline of the lungs and diaphragm may sit behind the highlight. Calm face, eyes closed. Aspect ratio 4:3 landscape. Check before use: certified teacher confirms the form against the steps in the notes block.
        % NEW-PIC (uncomment to show the new picture beside the original; comment the original only after review) \includegraphics[width=0.35\linewidth,keepaspectratio]{new/ycb1_threepartbreathing_new.png}
        \includegraphics[width=\linewidth,keepaspectratio]{ycb1_threepartbreathing}
				
		{\tiny (Ref: Beginner Yoga Flow)}	  
        \end{center}	
    \end{column}
\end{columns}
\end{frame}


%%%%%%%%%%%%%%%%%%%%%%%%%%%%%%%%%%%%%%%%%%%%%%%%%%%%%%%%%%%%%%%%%%%%%%%%%%%%%%%%%%
\begin{frame}[fragile]\frametitle{}
\begin{center}
{\Large प्राणायाम Pranayama}
\end{center}
\end{frame}

%%%%%%%%%%%%%%%%%%%%%%%%%%%%%%%%%%%%%%%%%%%%%%%%%%%%%%%%%%%
\begin{frame}[fragile]\frametitle{प्राणायाम Pranayama: Concept}

      \begin{itemize}
        \item \textbf{Pranayama}: regulating and lengthening the \textbf{breath}.
        \item \textbf{Prana} (life force) plus \textbf{ayama} (extension, control by stretching).
        \item The fourth of the eight limbs, after asana and before pratyahara.
        \item Traditionally said to \textbf{calm the mind}.
        \item Practised with or without \textbf{kumbhaka} (breath retention); the beginner practices are without it.
		\item  यम, नियम, आसन, प्राणायाम, प्रत्याहार, धारणा, ध्यान, तथा समाधि । प्राणायाम = प्राण + आयाम । इसका शाब्दिक अर्थ है -  प्राण या श्वसन को लम्बा करना  या फिर   जीवनी शक्ति  को लम्बा करना । प्राणायाम का अर्थ कुछ हद तक श्वास को नियंत्रित करना हो सकता है | परन्तु स्वास को कम करना नहीं होता है | प्राण या श्वास का आयाम या विस्तार ही प्राणायाम कहलाता है | यह प्राण-शक्ति का प्रवाह कर व्यक्ति को जीवन शक्ति प्रदान करता है।
		\item क्रिया: पूरक: श्वास घेणे , कुम्भक : रोखणे , रेचक : सोडणे 

      \end{itemize}

\end{frame}

%%%%%%%%%%%%%%%%%%%%%%%%%%%%%%%%%%%%%%%%%%%%%%%%%%%%%%%%%%%
\begin{frame}[fragile]\frametitle{Types of Pranayama}

      \begin{itemize}
        \item \textbf{अनुलोम विलोम  Anulom Vilom}: Alternate nostril \textbf{breathing}.
        \item \textbf{कपालभाती  Kapalabhati}: Skull \textbf{shining breath}; a cleansing kriya, not in the Board's pranayama list.
        \item \textbf{भ्रामरी Bhramari}: \textbf{Bee} breath.
        \item \textbf{उज्जायी Ujjayi}: \textbf{Victorious} breath.
        \item \textbf{शीतली Sheetali}: Cooling \textbf{breath}.
        \item घेरन्ड संहिता के अनुसार - सहित: सूर्यभेदश्च उज्जायी शीतली तथा । भस्त्रिका भरमारी मूर्च्छा केवली चाष्टकुम्भका: ।।
        \item घेरन्ड संहिता के अनुसार प्राणायाम के आठ भेद बताए गए हैं -सहित, सूर्यभेदी, उज्जायी, शीतली, भस्त्रिका, भ्रामरी, मूर्च्छा और केवली ।
        \item हठप्रदीपिका के अनुसार -सूर्यभेदनमुज़्जायी सीत्कारी शीतली तथा ।भस्त्रिका भ्रामरी मूर्च्छा प्लाविनीत्यष्ट कुम्भक: ।।
        \item हठप्रदीपिका के अनुसार प्राणायाम के आठ भेद निम्न हैं - सूर्यभेदन, उज्जायी, सीत्कारी, शीतली, भस्त्रिका, भ्रामरी, मूर्च्छा और प्लाविनी ये आठ प्रकार के कुम्भक (प्राणायाम ) होते हैं ।
		
      \end{itemize}

\end{frame}

%%%%%%%%%%%%%%%%%%%%%%%%%%%%%%%%%%%%%%%%%%%%%%%%%%%%%%%%%%%
\begin{frame}[fragile]\frametitle{Benefits of Pranayama}
      \begin{itemize}
        \item Traditionally credited with a calm mind and \textbf{peace of mind} (Guidebook).
        \item Links \textbf{breath} and \textbf{mind}; nadi shodhana lowers \textbf{stress} (Common Yoga Protocol).
        \item Bhramari calms the mind and prepares it for \textbf{concentration} (Common Yoga Protocol).
        % \item Supports \textbf{emotional stability}. % not in the notes

        \item These are traditional or general benefits, not medical claims.
      \end{itemize}

\end{frame}

%%%%%%%%%%%%%%%%%%%%%%%%%%%%%%%%%%%%%%%%%%%%%%%%%%%%%%%%%%%
\begin{frame}[fragile]\frametitle{अनुलोम विलोम Anulom Vilom (Alternate Nostril Breathing)}
\begin{columns}
    \begin{column}[T]{0.5\linewidth}
      \begin{itemize}
        \item Right hand in \textbf{nasagra mudra}, left hand in jnana mudra: inhale left, exhale right, inhale right, exhale left (one round).
        \item Traditionally said to purify the \textbf{nadis} and bring calm.
        \item Traditionally credited with a steady mind and lower \textbf{stress}.
        \item Slow, steady, never forced; equal inhale and exhale at first, later one to two.
        \item Five rounds (Common Yoga Protocol).
      \end{itemize}
    \end{column}
    \begin{column}[T]{0.5\linewidth}
        \begin{center}
        % FIG ycb1-fig-anulomvilom | now: images/ycb1_anulomvilom.png | final: books/YCB1/images/new/ycb1_anulomvilom_new.png | register: books/YCB1/docs/image_register.csv
        % PROMPT SP-FIG (common text: books/YCB1/docs/image_system_prompt.md) | Used in block ycb1-pd-pranayama, frame 'Anulom Vilom (Alternate Nostril Breathing)'. Pranayama: Anulom Vilom (Alternate Nostril Breathing). One adult practitioner (CS1), seated in a comfortable cross-legged pose, spine upright, eyes closed. Hands, fingers, face and tongue or lips exactly as in the notes block for this practice; mark the side of the breath with thin arrows only if the notes block gives a side. Front three-quarter view, head and shoulders large enough to show the hand and face. Aspect ratio 4:5 portrait. Check before use: certified teacher confirms the form against the steps in the notes block.
        % NEW-PIC (uncomment to show the new picture beside the original; comment the original only after review) \includegraphics[width=0.35\linewidth,keepaspectratio]{new/ycb1_anulomvilom_new.png}
        \includegraphics[width=\linewidth,keepaspectratio]{ycb1_anulomvilom}
				
		{\tiny (Ref: Certification  of Yoga Professionals Official Guidebook For Level I (Instructor))}	  
        \end{center}	
    \end{column}
\end{columns}
\end{frame}

%%%%%%%%%%%%%%%%%%%%%%%%%%%%%%%%%%%%%%%%%%%%%%%%%%%%%%%%%%%
\begin{frame}[fragile,shrink=6]\frametitle{कपालभाती  Kapalabhati (Skull Shining Breath)}
\begin{columns}
    \begin{column}[T]{0.5\linewidth}
      \begin{itemize}
        \item \textbf{Forceful exhalation} followed by \textbf{passive inhalation}.
        \item \textbf{Cleanses} the frontal region of the head and the lungs (Guidebook).
        \item Forceful exhalation by the abdominal muscles; chest and shoulders stay still.
        \item Thirty rapid breaths, then a deep breath; three rounds (beginners up to three rounds of twenty).
        \item Avoid with cardiac conditions, giddiness, \textbf{high blood pressure}, vertigo, chronic nose bleeding, epilepsy, migraine, stroke, hernia or gastric ulcer.
      \end{itemize}
    \end{column}
    \begin{column}[T]{0.5\linewidth}
        \begin{center}
        % FIG ycb1-fig-kapalbhati | now: images/ycb1_kapalbhati.png | final: books/YCB1/images/new/ycb1_kapalbhati_new.png | register: books/YCB1/docs/image_register.csv
        % PROMPT SP-FIG (common text: books/YCB1/docs/image_system_prompt.md) | Used in block ycb1-pd-cleansing-techniques, frame 'Kapalabhati (Skull Shining Breath); Kapalabhati'. Practice: Kapalabhati (Skull Shining Breath); Kapalabhati. One adult practitioner (CS1), seated, front three-quarter view, in the posture of the notes block. Show the key action in a small inset: the abdomen drawn in sharply for kapalabhati, the throat narrowed softly for ujjayi, the steady gaze on a candle flame at eye level for trataka, the tongue position for the mukh mudra. Use only the inset that matches Kapalabhati (Skull Shining Breath); Kapalabhati. Aspect ratio 4:5 portrait. Check before use: certified teacher confirms the form against the steps in the notes block.
        % NEW-PIC (uncomment to show the new picture beside the original; comment the original only after review) \includegraphics[width=0.35\linewidth,keepaspectratio]{new/ycb1_kapalbhati_new.png}
        \includegraphics[width=0.8\linewidth,keepaspectratio]{ycb1_kapalbhati}
				
		{\tiny (Ref: Certification  of Yoga Professionals Official Guidebook For Level I (Instructor))}	 
        \end{center}	
    \end{column}
\end{columns}
\end{frame}

%%%%%%%%%%%%%%%%%%%%%%%%%%%%%%%%%%%%%%%%%%%%%%%%%%%%%%%%%%%
\begin{frame}[fragile]\frametitle{भ्रामरी  Bhramari (Bee Breath)}
\begin{columns}
    \begin{column}[T]{0.5\linewidth}
      \begin{itemize}
        \item \textbf{Inhale} deeply and \textbf{exhale} with a humming sound.
        \item Traditionally credited with easing \textbf{stress} and calming the mind.
        \item A preparation for \textbf{concentration} and meditation.
        \item Five rounds (Common Yoga Protocol); avoid with nose or ear infections.
        \item The second form uses \textbf{shanmukhi mudra} (eyes, mouth and ears closed with the fingers).
      \end{itemize}
    \end{column}
    \begin{column}[T]{0.5\linewidth}
        \begin{center}
        % FIG ycb1-fig-bhramari | now: images/ycb1_bhramari.png | final: books/YCB1/images/new/ycb1_bhramari_new.png | register: books/YCB1/docs/image_register.csv
        % PROMPT SP-FIG (common text: books/YCB1/docs/image_system_prompt.md) | Used in block ycb1-pd-pranayama, frame 'Bhramari (Bee Breath)'. Pranayama: Bhramari (Bee Breath). One adult practitioner (CS1), seated in a comfortable cross-legged pose, spine upright, eyes closed. Hands, fingers, face and tongue or lips exactly as in the notes block for this practice; mark the side of the breath with thin arrows only if the notes block gives a side. Front three-quarter view, head and shoulders large enough to show the hand and face. Aspect ratio 4:5 portrait. Check before use: certified teacher confirms the form against the steps in the notes block.
        % NEW-PIC (uncomment to show the new picture beside the original; comment the original only after review) \includegraphics[width=0.35\linewidth,keepaspectratio]{new/ycb1_bhramari_new.png}
        \includegraphics[width=\linewidth,keepaspectratio]{ycb1_bhramari}
				
		{\tiny (Ref: Certification  of Yoga Professionals Official Guidebook For Level I (Instructor))}	 
        \end{center}	
    \end{column}
\end{columns}
\end{frame}

%%%%%%%%%%%%%%%%%%%%%%%%%%%%%%%%%%%%%%%%%%%%%%%%%%%%%%%%%%%
\begin{frame}[fragile]\frametitle{उज्जायी Ujjayi (Victorious Breath)}
\begin{columns}
    \begin{column}[T]{0.5\linewidth}
      \begin{itemize}
        \item \textbf{Inhale} and \textbf{exhale} with a slight constriction of the throat.
        \item Produces a \textbf{soothing} and \textbf{hissing} sound.
        \item Can be done sitting, standing or lying (Guidebook).
        \item The throat is only partly closed; never strain.
        \item Soft, even sound on inhale and exhale.
      \end{itemize}
    \end{column}
    \begin{column}[T]{0.5\linewidth}
        \begin{center}
        % FIG ycb1-fig-ujjayi | now: images/ycb1_ujjayi.png | final: books/YCB1/images/new/ycb1_ujjayi_new.png | register: books/YCB1/docs/image_register.csv
        % PROMPT SP-FIG (common text: books/YCB1/docs/image_system_prompt.md) | Used in block ycb1-pd-pranayama, frame 'Ujjayi (Victorious Breath)'. Practice: Ujjayi (Victorious Breath). One adult practitioner (CS1), seated, front three-quarter view, in the posture of the notes block. Show the key action in a small inset: the abdomen drawn in sharply for kapalabhati, the throat narrowed softly for ujjayi, the steady gaze on a candle flame at eye level for trataka, the tongue position for the mukh mudra. Use only the inset that matches Ujjayi (Victorious Breath). Aspect ratio 4:5 portrait. Check before use: certified teacher confirms the form against the steps in the notes block.
        % NEW-PIC (uncomment to show the new picture beside the original; comment the original only after review) \includegraphics[width=0.35\linewidth,keepaspectratio]{new/ycb1_ujjayi_new.png}
        \includegraphics[width=\linewidth,keepaspectratio]{ycb1_ujjayi}
				
		{\tiny (Ref: Certification  of Yoga Professionals Official Guidebook For Level I (Instructor))}	 
        \end{center}	
    \end{column}
\end{columns}
\end{frame}

%%%%%%%%%%%%%%%%%%%%%%%%%%%%%%%%%%%%%%%%%%%%%%%%%%%%%%%%%%%
\begin{frame}[fragile]\frametitle{शीतली Sheetali (Cooling Breath)}
\begin{columns}
    \begin{column}[T]{0.5\linewidth}
      \begin{itemize}
        \item \textbf{Inhale} through the rolled tongue (sheetkari: through the teeth if the tongue cannot roll).
        \item \textbf{Exhale} through the nose.
        \item Cools the \textbf{body} and \textbf{mind}.
        \item Calms the \textbf{mind}; never force; stop if dizzy.
        \item Avoid with cold, cough, asthma, tonsillitis or low blood pressure, and in cold weather.
      \end{itemize}
    \end{column}
    \begin{column}[T]{0.5\linewidth}
        \begin{center}
        % FIG ycb1-fig-sitali | now: images/ycb1_sitali.png | final: books/YCB1/images/new/ycb1_sitali_new.png | register: books/YCB1/docs/image_register.csv
        % PROMPT SP-FIG (common text: books/YCB1/docs/image_system_prompt.md) | Used in block ycb1-pd-pranayama, frame 'Sitali (Cooling Breath)'. Pranayama: Sitali (Cooling Breath). One adult practitioner (CS1), seated in a comfortable cross-legged pose, spine upright, eyes closed. Hands, fingers, face and tongue or lips exactly as in the notes block for this practice; mark the side of the breath with thin arrows only if the notes block gives a side. Front three-quarter view, head and shoulders large enough to show the hand and face. Aspect ratio 4:5 portrait. Check before use: certified teacher confirms the form against the steps in the notes block.
        % NEW-PIC (uncomment to show the new picture beside the original; comment the original only after review) \includegraphics[width=0.35\linewidth,keepaspectratio]{new/ycb1_sitali_new.png}
        \includegraphics[width=\linewidth,keepaspectratio]{ycb1_sitali}
				
		{\tiny (Ref: Certification  of Yoga Professionals Official Guidebook For Level I (Instructor))}	 
        \end{center}	
    \end{column}
\end{columns}
\end{frame}


%%%%%%%%%%%%%%%%%%%%%%%%%%%%%%%%%%%%%%%%%%%%%%%%%%%%%%%%%%%%%%%%%%%%%%%%%%%%%%%%%%
\begin{frame}[fragile]\frametitle{}
\begin{center}
{\Large Understanding of मुद्रा  Mudra}
\end{center}
\end{frame}

%%%%%%%%%%%%%%%%%%%%%%%%%%%%%%%%%%%%%%%%%%%%%%%%%%%%%%%%%%%
\begin{frame}[fragile]\frametitle{Understanding of मुद्रा  Mudra: Overview}

      \begin{itemize}
        \item \textbf{मुद्रा  Mudra}: a \textbf{seal} or \textbf{gesture}; here, of the \textbf{hands}.
        \item The word is Sanskrit for \textbf{seal} or \textbf{gesture}.
        \item In tradition, said to \textbf{channel} and \textbf{direct energy}.  प्राण शक्ती
        \item Used in \textbf{meditation} and in \textbf{प्राणायाम  pranayama}.
        % \item Integrates with \textbf{आसन   asanas} and \textbf{प्राणायाम  pranayama}. % asanas: not in the notes (pending block)
		\item A traditional idea: energy is said to leave through the fingers, and the mudra turns it back into the body.
      \end{itemize}
\end{frame}

%%%%%%%%%%%%%%%%%%%%%%%%%%%%%%%%%%%%%%%%%%%%%%%%%%%%%%%%%%%
\begin{frame}[fragile]\frametitle{Types of Mudras}

      \begin{itemize}
        \item The Board syllabus asks only for \textbf{Hasta Mudras (हस्त मुद्रा)}: hand \textbf{gestures}.
        \item The Bihar School groups mudras as hasta, mana (head), kaya, bandha and adhara (Satyananda).
        \item Bandha and mudra are separate topics in the syllabus.
      \end{itemize}

\end{frame}


%%%%%%%%%%%%%%%%%%%%%%%%%%%%%%%%%%%%%%%%%%%%%%%%%%%%%%%%%%%
\begin{frame}[fragile,shrink=6]\frametitle{Hasta Mudras (Hand Gestures)}
\begin{columns}
    \begin{column}[T]{0.64\linewidth}
      \footnotesize
      \begin{itemize}
        \item \textbf{Gyan Mudra (ज्ञान मुद्रा)}: \textbf{Knowledge} gesture. First finger tip touching the thumb tip in a circle, other three fingers straight, relaxed and together; palms up on the thighs in the Common Yoga Protocol (some texts use palms down).
        \item \textbf{Chin Mudra (चिन मुद्रा)}: \textbf{Consciousness} gesture. Same finger contact as jnana. The Board syllabus lists it separately and sources differ on palm direction, so confirm with a certified teacher.
        \item \textbf{Chinmay Mudra (चिन्मय  मुद्रा)}: \textbf{Manifest consciousness} gesture. Jnana hold, then the other three fingers folded in towards the palm; palms up or down (Ref: Niranjananand, Pran and Pranayam).	
		\item \textbf{Brahma Mudra (ब्रह्म   मुद्रा)}: \textbf{Divine} gesture. Thumbs folded inside across the palms, fingers folded over them to form fists; backs of the fists rest on the thighs, knuckles touching, close to the body at pelvic level (Ref: Niranjananand, Pran and Pranayam).
		\item \textbf{Adi Mudra (आदी   मुद्रा)}: Thumb folded into the palm and fingers curled over it to make a fist; hands on the knees, palms up or down (Ref: Niranjananand, Pran and Pranayam).
		\item \textbf{Dhyan Mudra (ध्यान    मुद्रा)}: Right hand on the left, palms up, resting in the lap. Satyananda names this form bhairava mudra; the Board must confirm what it calls dhyana mudra. (The Common Yoga Protocol meditates in jnana mudra.)
		\item \textbf{Nasika/Vishnu Mudra (नासिका / विष्णू     मुद्रा)}: First two fingers folded, ring and little fingers by the left nostril, thumb by the right; called nasagra mudra in the Common Yoga Protocol, used in Anulom Vilom.
        % \item \textbf{Anjali Mudra (अंजलि मुद्रा)}: \textbf{Salutation} gesture.
        % \item \textbf{Apan Mudra (अपान मुद्रा)}: \textbf{Cleansing} gesture.
        % \item \textbf{Shuni Mudra (शुणि मुद्रा)}: \textbf{Patience} gesture.
      \end{itemize}
    \end{column}
    \begin{column}[T]{0.33\linewidth}
        \begin{center}
        % FIG ycb1-fig-hastamudra | now: images/ycb1_hastamudra.jpg | final: books/YCB1/images/new/ycb1_hastamudra_new.png | register: books/YCB1/docs/image_register.csv
        % PROMPT SP-FIG (common text: books/YCB1/docs/image_system_prompt.md) | Used in block ycb1-pd-mudra, frame 'Hasta Mudras (Hand Gestures)'. Mudra practice: Hasta Mudras (Hand Gestures). Close-up of the hand or hands in the gesture named in the notes block, one hand per panel if more than one gesture is shown, same skin tone in every panel, five fingers per hand, fingertips and thumb contact clearly drawn. Small inset of the seated practitioner (CS1) showing where the hands rest. Aspect ratio 4:3. Check before use: certified teacher confirms the form against the steps in the notes block.
        % NEW-PIC (uncomment to show the new picture beside the original; comment the original only after review) \includegraphics[width=0.35\linewidth,keepaspectratio]{new/ycb1_hastamudra_new.png}
        \includegraphics[width=\linewidth,keepaspectratio]{ycb1_hastamudra}
				
		{\tiny (Ref: Himalayan Yoga Academy)}	 
        \end{center}	
    \end{column}
\end{columns}
\end{frame}


%%%%%%%%%%%%%%%%%%%%%%%%%%%%%%%%%%%%%%%%%%%%%%%%%%%%%%%%%%%
\begin{frame}[fragile]\frametitle{Kaya Mudras (Body Gestures)}
\begin{columns}
    \begin{column}[T]{0.5\linewidth}
      \begin{itemize}
        \item \textbf{Kaya mudras}: a school grouping of body-based practices; not a Board syllabus topic (the syllabus lists hasta mudras).
        \item \textbf{Viparita Karani (विपरीत करणी)}: an asana, with a caution for neck and blood pressure.
        \item \textbf{Sarvangasana (सर्वांगासन)}: an asana, with a caution for neck and blood pressure.
        \item \textbf{Adho Mukha Svanasana (अधोमुख श्वानासन)}: an asana, not a mudra of the syllabus.
        \item Learn these in the asana block.
      \end{itemize}
    \end{column}
    \begin{column}[T]{0.5\linewidth}
        \begin{center}
        % FIG ycb1-fig-kayamudra | now: images/ycb1_kayamudra.jpg | final: books/YCB1/images/new/ycb1_kayamudra_new.png | register: books/YCB1/docs/image_register.csv
        % PROMPT SP-FIG (common text: books/YCB1/docs/image_system_prompt.md) | Used in block ycb1-pd-mudra, frame 'Kaya Mudras (Body Gestures)'. Mudra practice: Kaya Mudras (Body Gestures). Close-up of the hand or hands in the gesture named in the notes block, one hand per panel if more than one gesture is shown, same skin tone in every panel, five fingers per hand, fingertips and thumb contact clearly drawn. Small inset of the seated practitioner (CS1) showing where the hands rest. Aspect ratio 4:3. Check before use: certified teacher confirms the form against the steps in the notes block.
        % NEW-PIC (uncomment to show the new picture beside the original; comment the original only after review) \includegraphics[width=0.35\linewidth,keepaspectratio]{new/ycb1_kayamudra_new.png}
        \includegraphics[width=\linewidth,keepaspectratio]{ycb1_kayamudra}
				
		{\tiny (Ref: Prana Sutra)}	 
        \end{center}	
    \end{column}
\end{columns}
\end{frame}


%%%%%%%%%%%%%%%%%%%%%%%%%%%%%%%%%%%%%%%%%%%%%%%%%%%%%%%%%%%
% Commented out 2026-11-27 (outside the syllabus; pass 4 resolution)
% \begin{frame}[fragile]\frametitle{Mukh Mudras (Facial Gestures)}
% \begin{columns}
%     \begin{column}[T]{0.5\linewidth}
%       \begin{itemize}
%         \item \textbf{Khechari Mudra (खेचरी मुद्रा)}: \textbf{Tongue} gesture.
%         \item \textbf{Bhrumadhya Mudra (भ्रूमध्य मुद्रा)}: \textbf{Eyebrow} gesture.
%         \item \textbf{Shambhavi Mudra (शाम्भवी  मुद्रा)}: \textbf{Eyebrow center} gaze.
%         \item Traditionally used to turn the attention inward.
%         \item Integrates with \textbf{meditation} and \textbf{pranayama}.
%       \end{itemize}
%     \end{column}
%     \begin{column}[T]{0.5\linewidth}
%         \begin{center}
%         % FIG ycb1-fig-khechari | now: images/ycb1_khechari.png | final: books/YCB1/images/new/ycb1_khechari_new.png | register: books/YCB1/docs/image_register.csv
%         % PROMPT SP-FIG (common text: books/YCB1/docs/image_system_prompt.md) | Used in block ycb1-pd-pranayama, frame 'Mukh Mudras (Facial Gestures)'. Practice: Mukh Mudras (Facial Gestures). One adult practitioner (CS1), seated, front three-quarter view, in the posture of the notes block. Show the key action in a small inset: the abdomen drawn in sharply for kapalabhati, the throat narrowed softly for ujjayi, the steady gaze on a candle flame at eye level for trataka, the tongue position for the mukh mudra. Use only the inset that matches Mukh Mudras (Facial Gestures). Aspect ratio 4:5 portrait. Check before use: certified teacher confirms the form against the steps in the notes block.
%         % NEW-PIC (uncomment to show the new picture beside the original; comment the original only after review) \includegraphics[width=0.35\linewidth,keepaspectratio]{new/ycb1_khechari_new.png}
%         \includegraphics[width=\linewidth,keepaspectratio]{ycb1_khechari}
				
% 		{\tiny (Ref: Certification of Yoga Professionals Official Guidebook For Level I (Instructor))}	 
%         \end{center}	
%     \end{column}
% \end{columns}
% \end{frame}


%%%%%%%%%%%%%%%%%%%%%%%%%%%%%%%%%%%%%%%%%%%%%%%%%%%%%%%%%%%
\begin{frame}[fragile,shrink=6]\frametitle{बंध  Bandhas (Internal Locks)}\footnotesize

% बंध मुद्रा शरीर की कुछ ऐसी अवस्था हैं जिनके द्वारा कुंडलिनी सफलतापूर्वक जाग्रत की जा सकती है। घेरंड संहिता में २५ मुद्राओं एवं महत्वपूर्ण हैं:
% (१) मूलबंध, (२) जालंधरबंध, (३) उड्डीयानबंध, (४) महामुद्रा, (५) महाबंध, (६) महावेध
% (७) योगमुद्रा, (८) विपरीतकरणीमुद्रा, (९) खेचरीमुद्रा, (१०) वज्रिणीमुद्रा, (११) शक्तिचालिनीमुद्रा, (१२) योनिमुद्रा।

\begin{columns}
    \begin{column}[T]{0.5\linewidth}
      \begin{itemize}
        \item \textbf{Mula Bandha (मूल बन्ध)}: Root \textbf{lock}: the perineum is pressed and the anus contracted (Hatha Pradipika).
        \item \textbf{Uddiyana Bandha (उड्डीयान बन्ध)}: Abdominal \textbf{lock}: after a full exhalation the abdomen is drawn in and up.
        \item \textbf{Jalandhara Bandha (जालंधर बन्ध)}: Throat \textbf{lock}: chin pressed firmly on the chest (Hatha Pradipika).
        \item Done with pranayama, after the asanas; traditionally said to direct \textbf{prana}.
        \item \textbf{Uddiyana}: empty stomach only. Avoid with colitis, ulcer, hernia, high blood pressure, heart disease, glaucoma, raised intracranial pressure, pregnancy, menstruation.
        \item \textbf{Jalandhara}: avoid with cervical spondylosis, raised intracranial pressure, vertigo, high blood pressure or heart disease.
        \item Learn all three from a \textbf{qualified teacher}; never force. Mula only under an experienced teacher.
      \end{itemize}
    \end{column}
    \begin{column}[T]{0.5\linewidth}
        \begin{center}
        % FIG ycb1-fig-bandhas | now: images/ycb1_bandhas.jpg | final: books/YCB1/images/new/ycb1_bandhas_new.pdf | register: books/YCB1/docs/image_register.csv
        % PROMPT SP-DIAGRAM (common text: books/YCB1/docs/image_system_prompt.md) | Used in block ycb1-pd-bandha, frame 'Bandhas (Internal Locks)'. Concept diagram: Bandhas (Internal Locks). Boxes, circles and arrows only, drawn from the structure given in the notes block (the names and their order are listed there); Sanskrit terms in Roman letters with the English meaning in brackets, nothing else added. Flat colours that still read in grayscale. Best redrawn later in TikZ. Aspect ratio 4:3. Check before use: author and the Sanskrit check in the notes.
        % NEW-PIC approved by the author (2026-11-27); the original is commented out below it
        \includegraphics[width=\linewidth,keepaspectratio]{new/ycb1_bandhas_new.pdf}
        % original, replaced (author approved 2026-11-27): \includegraphics[width=\linewidth,keepaspectratio]{ycb1_bandhas}
				
		% credit of the replaced original: {\tiny (Ref: Arogya Yoga School)}
		{\tiny (Line diagram drawn for this deck.)}
        \end{center}	
    \end{column}
\end{columns}
\end{frame}


%%%%%%%%%%%%%%%%%%%%%%%%%%%%%%%%%%%%%%%%%%%%%%%%%%%%%%%%%%%
% Commented out 2026-11-27 (outside the syllabus; pass 4 resolution)
% \begin{frame}[fragile]\frametitle{Chakra Mudras (Energy Center Gestures)}
% \begin{columns}
%     \begin{column}[T]{0.5\linewidth}
%       \begin{itemize}
%         \item \textbf{Root Chakra Mudra (मूलाधार चक्र मुद्रा)}: \textbf{Grounding} gesture.
%         \item \textbf{Heart Chakra Mudra (अनाहत चक्र मुद्रा)}: \textbf{Love} gesture.
%         \item \textbf{Third Eye Chakra Mudra (आज्ञा चक्र मुद्रा)}: \textbf{Intuition} gesture.
%         \item In tradition, linked to the \textbf{energy centers} and \textbf{meditation}.
%         \item A school classification; forms vary between teachers.
%       \end{itemize}
%     \end{column}
%     \begin{column}[T]{0.5\linewidth}
%         \begin{center}
%         % FIG ycb1-fig-chakras | now: images/ycb1_chakras.jpg | final: books/YCB1/images/new/ycb1_chakras_new.pdf | register: books/YCB1/docs/image_register.csv
%         % PROMPT SP-DIAGRAM (common text: books/YCB1/docs/image_system_prompt.md) | Used in block ycb1-th-principles-of-yoga, frame 'Chakra Mudras (Energy Center Gestures); Chakras'. Concept diagram: Chakra Mudras (Energy Center Gestures); Chakras. Boxes, circles and arrows only, drawn from the structure given in the notes block (the names and their order are listed there); Sanskrit terms in Roman letters with the English meaning in brackets, nothing else added. Flat colours that still read in grayscale. Best redrawn later in TikZ. Aspect ratio 4:3. Check before use: author and the Sanskrit check in the notes.
%         % NEW-PIC (uncomment to show the new picture beside the original; comment the original only after review) \includegraphics[width=0.35\linewidth,keepaspectratio]{new/ycb1_chakras_new.pdf}
%         \includegraphics[width=\linewidth,keepaspectratio]{ycb1_chakras}
				
% 		{\tiny (Ref: Raja Yoga Rishikesh)}	 
%         \end{center}	
%     \end{column}
% \end{columns}
% \end{frame}



%%%%%%%%%%%%%%%%%%%%%%%%%%%%%%%%%%%%%%%%%%%%%%%%%%%%%%%%%%%%%%%%%%%%%%%%%%%%%%%%%%
\begin{frame}[fragile]\frametitle{}
\begin{center}
{\Large Practices leading to Meditation and ध्यान साधना  Dhyana Sadhana }
\end{center}
\end{frame}

%%%%%%%%%%%%%%%%%%%%%%%%%%%%%%%%%%%%%%%%%%%%%%%%%%%%%%%%%%%
\begin{frame}[fragile]\frametitle{Practices Leading to Meditation and ध्यान साधना  Dhyana Sadhana}

      \begin{itemize}
        \item One way to prepare (a teaching sequence, not the Board order; the syllabus lists dharana and dhyana with mudras, mantra and japa):
        \item \textbf{Mindfulness}: Develop \textbf{awareness} of \textbf{thoughts} and \textbf{emotions}.
        \item \textbf{Breathing Techniques}: Practice \textbf{Pranayama} to calm the \textbf{mind}.
        \item \textbf{Asanas}: Perform \textbf{stabilizing poses} to prepare for \textbf{meditation}.
        \item \textbf{Concentration Exercises}: Engage in \textbf{focusing} techniques to enhance \textbf{mental clarity}.
        \item \textbf{Visualization}: Use \textbf{guided imagery} to support \textbf{meditative focus}.
      \end{itemize}
 
\end{frame}








